
\dictentry{Ablation}{The removal of surface material through vaporization, melting, or sublimation to dissipate intense aerodynamic heating during atmospheric entry. Ablation absorbs heat energy through phase change and mass loss, protecting the underlying structure from thermal damage. The rate and depth of ablation depend on the material properties, heat flux, and duration of exposure.}{catAD}
\dictentry{Ablative TPS}{A thermal protection system that dissipates heat through the controlled ablation (surface removal) of a sacrificial material, carrying heat away from the underlying structure. Ablative TPS is used for high-heat-flux applications such as ICBM re-entry vehicles and planetary probes.}{catAD}
\dictentry{Accelerometer}{A transducer that converts mechanical acceleration into an electrical signal, used to measure vibration, shock, and dynamic loads on aircraft structures and components. Accelerometers operate on piezoelectric, capacitive, or MEMS principles and are essential for flight test instrumentation and ground vibration testing. They provide data for structural dynamic analysis, flutter clearance, and ride quality assessment.}{catAD}
\dictentry{Acoustic Testing}{A specialized wind tunnel testing methodology focused on measuring noise generated by aerodynamic flows. It requires anechoic or treated test sections with low background noise levels to accurately capture airframe, engine, and propeller noise signatures.}{catAD}
\dictentry{Active Flow Control}{A method of modifying aerodynamic flow fields using energy input such as blowing, suction, plasma actuators, or synthetic jets that can be actively modulated in real time. Active flow control can delay separation, enhance mixing, reduce drag, and augment lift beyond static geometric limits. It requires power, sensors, and control systems but offers adaptability to changing flight conditions. Applications include separation control on airfoils, noise reduction in jet exhaust, and cavity flow management.}{catAD}

\dictentry{Adverse Yaw}{The yawing tendency of an aircraft in the direction opposite to the intended roll, caused by the increased induced drag on the wing that is generating more lift during a turn. Adverse yaw is most pronounced at low speeds and high angles of attack, and is counteracted through differential aileron deflection, wing sweep, or coordinated rudder input.}{catAD}
\dictentry{Aeroacoustics}{The branch of acoustics that studies noise generation by turbulent fluid motion or aerodynamic forces interacting with surfaces. It encompasses jet noise, airframe noise, propeller and rotor noise, and fan noise, and is critical for aircraft certification and community noise reduction.}{catAD}
\dictentry{Aerodynamic Center}{The point on an airfoil or wing about which the pitching moment coefficient remains approximately constant regardless of angle of attack. For subsonic flow it typically lies near the quarter-chord point, while at supersonic speeds it shifts rearward to roughly the mid-chord. Knowing the aerodynamic center is essential for longitudinal stability analysis and trim calculations.}{catAD}
\dictentry{Aerodynamic Force}{The resultant force exerted on a body moving through the atmosphere, resolved into three components: lift perpendicular to the freestream, drag opposing the motion, and side force perpendicular to both. These forces arise from pressure distributions and shear stresses on the body surface. Aerodynamic force magnitudes depend on airspeed, air $\rho$, reference area, and the body's shape and orientation.}{catAD}
\dictentry{Aerodynamic Heating}{The process by which kinetic energy of high-speed airflow is converted to thermal energy on a vehicle's surface through viscous dissipation and compression. It becomes significant at transonic speeds and increases dramatically in the hypersonic regime, requiring thermal protection systems for vehicle survival. The heating rate depends on velocity cubed, air $\rho$, and the local surface geometry.}{catAD}
\dictentry{Aerodynamic Twist}{The variation in effective angle of attack along a wing's span due to factors such as washout, sweep, or fuselage upwash. Aerodynamic twist influences the spanwise lift distribution and is deliberately designed into wings to control stall progression and improve lateral stability. Unlike geometric twist, aerodynamic twist can change with flight condition.}{catAD}
\dictentry{Aerodynamics}{The branch of fluid mechanics concerned with the motion of air and the forces acting on bodies immersed in it. It encompasses the study of lift, drag, flow separation, boundary layers, and compressibility effects, and is fundamental to the design of aircraft, missiles, and other flight vehicles. Modern aerodynamics combines theoretical analysis, computational simulation, and wind tunnel experimentation.}{catAD}
\dictentry{Aeroelasticity}{The study of the coupled interaction between aerodynamic forces, elastic structural deformation, and inertial effects. Aeroelastic phenomena include flutter, divergence, and control reversal, which can lead to catastrophic structural failure if not properly addressed during design. It is critical for lightweight, flexible structures such as long-span wings, helicopter rotors, and high-aspect-ratio UAVs.}{catAD}
\dictentry{Aeroshell}{The protective outer shell of a spacecraft that provides aerodynamic deceleration and thermal protection during atmospheric entry. An aeroshell typically consists of a heat shield on the forward face and a backshell covering the aft side, enclosing the payload.}{catAD}
\dictentry{Aerothermal Heating}{The combined heating effect on a vehicle resulting from both convective heat transfer from the hot boundary layer and radiative emission from the shock layer during high-speed atmospheric flight. Aerothermal heating is particularly severe during hypersonic flight and atmospheric re-entry, where temperatures can exceed thousands of degrees. Accurate prediction of aerothermal heating is essential for thermal protection system design.}{catAD}
\dictentry{Agility}{The capability of an aircraft to rapidly change its attitude, including roll rate, pitch rate, and yaw rate, in response to pilot commands or autopilot inputs. High agility is essential for fighter aircraft engaged in air combat maneuvering and for missile intercept trajectories. It depends on control surface sizing, thrust-to-weight ratio, and the aircraft's moment of inertia about each axis.}{catAD}
\dictentry{Aileron}{A hinged control surface mounted on the trailing edge of a wing, typically outboard of the wing root, that controls the aircraft's rolling motion. Deflecting the aileron upward on one wing and downward on the other creates a differential lift that produces a rolling moment about the longitudinal axis. Ailerons are one of the three primary flight control surfaces along with the elevator and rudder.}{catAD}
\dictentry{Airbrake}{A deployable device that increases aerodynamic drag to decelerate an aircraft during flight, landing, or descent. Airbrakes can be mounted on the fuselage, wings, or tail and are used on military aircraft for rapid deceleration during combat maneuvers and on gliders for energy management during approach. They provide controlled drag without significantly altering lift or flight path.}{catAD}
\dictentry{Airfoil}{A two-dimensional cross-sectional shape specifically designed to generate lift or provide flow control when exposed to an airflow. Common airfoil families include NACA four-digit and five-digit series, supercritical, and laminar flow profiles, each optimized for specific speed ranges and applications. The performance of an airfoil is characterized by its lift coefficient, drag coefficient, and moment coefficient as functions of angle of attack.}{catAD}
\dictentry{Airframe Noise}{The aerodynamic noise produced by the airframe of an aircraft during flight, including noise from landing gear, flaps, slats, and wing trailing edges. As engines become quieter, airframe noise has become a dominant contributor to overall aircraft noise during approach and landing.}{catAD}
\dictentry{Angle of Attack}{The angle between the chord line of an airfoil or wing reference line and the oncoming relative wind vector, measured in the plane perpendicular to the span. Angle of attack is the primary variable controlling lift generation, with lift increasing approximately linearly until the critical angle is reached and the wing stalls. It is distinct from the pitch attitude and depends on the aircraft's flight path and sideslip.}{catAD}
\dictentry{Angle of Incidence}{The angle between the wing chord line and the aircraft's longitudinal (fuselage reference) axis, typically set during design and not adjustable in flight. A positive angle of incidence positions the wing at a slightly higher angle than the fuselage, optimizing cruise lift while keeping the fuselage level. This fixed geometric angle affects trim, drag, and the aircraft's attitude in level flight.}{catAD}
\dictentry{Anhedral}{A downward angular inclination of the wing from root to tip, measured relative to the lateral axis. Anhedral reduces lateral stability and is used on high-wing transport aircraft and some military fighters to counteract the excessive Dutch roll tendency produced by the high-mounted wing. It improves maneuverability and roll responsiveness at the expense of inherent lateral-directional stability.}{catAD}
\dictentry{Arc-Heated Wind Tunnel}{A hypersonic wind tunnel that uses a high-current electric arc to heat the test gas to very high temperatures, simulating the stagnation conditions encountered during atmospheric re-entry. Arc-heated facilities can produce stagnation temperatures exceeding 5000 K and are used for thermal protection material testing and aerothermal research. The arc heats the gas upstream of a converging-diverging nozzle that accelerates the flow to hypersonic speeds.}{catAD}
\dictentry{Area Rule}{A transonic aerodynamic design principle developed by Richard Whitcomb stating that the cross-sectional area of an aircraft should vary smoothly along its length to minimize wave drag at transonic speeds. The area rule led to the characteristic 'Coke bottle' fuselage shape of transonic aircraft and is critical for optimizing aerodynamic performance near Mach 1.}{catAD}
\dictentry{Aspect Ratio}{The ratio of a wing's span squared to its planform area, or equivalently the span divided by the mean aerodynamic chord. A high aspect ratio reduces induced drag and improves aerodynamic efficiency, which is why gliders and long-range cruise aircraft have long, narrow wings. Conversely, low aspect ratio wings favor high-speed performance and structural compactness.}{catAD}
\dictentry{Atmospheric Re-entry}{The phase of a spacecraft's return trajectory from space through a planetary atmosphere, characterized by hypersonic velocities, intense aerodynamic heating, and high deceleration loads. Atmospheric re-entry requires thermal protection systems and careful trajectory design to ensure survival.}{catAD}
\dictentry{Attached Flow}{Airflow that remains smoothly attached to and follows the contour of a surface without separating, producing well-defined pressure distributions and predictable aerodynamic forces. Attached flow is essential for efficient lift generation and is the design condition for airfoils, wings, and control surfaces. Flow attachment can be disrupted by excessive angle of attack, surface roughness, or adverse pressure gradients.}{catAD}
\dictentry{Attached Shock}{A shock wave that originates directly from the leading edge or sharp corner of a body in supersonic flow, remaining connected to the point of generation. Attached shocks form on sharp-nosed bodies at moderate angles of attack and produce weaker pressure jumps than detached bow shocks. The shock angle depends on freestream Mach number and body geometry.}{catAD}
\dictentry{Autorotation}{A state in which a helicopter rotor is driven solely by the upward flow of air through the rotor disc, without engine power, allowing a controlled descent and landing. Autorotation occurs when the aircraft descends at a sufficient rate such that the relative wind from below drives the blades. Pilots practice autorotation as an emergency procedure for engine failure, using stored kinetic energy to cushion the final landing flare.}{catAD}
\dictentry{Balance}{A multi-component force measurement device installed in a wind tunnel model to measure the aerodynamic forces (lift, drag, side force) and moments (pitching, rolling, yawing) acting on the model during testing. Balances use strain gauges or piezoelectric sensors arranged in Wheatstone bridge circuits to independently resolve each force and moment component. Accuracy and calibration of the balance are critical for obtaining reliable wind tunnel data.}{catAD}
\dictentry{Ballistic Re-entry}{A non-lifting re-entry trajectory in which the vehicle follows a purely gravitational path with no aerodynamic lift. Ballistic re-entry produces high peak deceleration and heating rates but is simpler and was used by early crewed capsules such as Mercury and Apollo.}{catAD}
\dictentry{Base Drag}{The component of pressure drag caused by the low-pressure wake region that forms behind a blunt base or truncated aft body where the flow separates. Base drag can account for a significant fraction of total drag on projectiles, rockets, and aircraft with blunt afterbodies. It is reduced by boat-tailing the aft section, using base bleed, or employing boat-tail nozzle designs.}{catAD}
\dictentry{Belly Flap}{A flap or control surface located on the underside (belly) of a spacecraft or lifting body, used to generate drag or modulate lift during atmospheric re-entry and descent. Belly flaps provide trim control and help manage the vehicle's trajectory and heating profile without requiring conventional wings. They were used on capsules like Apollo and are employed on modern re-entry vehicles for precision landing.}{catAD}
\dictentry{Bernoulli Equation}{A fundamental relation in inviscid, incompressible flow stating that the sum of static pressure, dynamic pressure, and hydrostatic pressure remains constant along a streamline. It is derived from the conservation of energy for a frictionless fluid and relates pressure changes to velocity changes. The equation underpins much of classical aerodynamics, including the explanation of lift generation on airfoils.}{catAD}
\dictentry{Best Glide}{The airspeed and corresponding lift-to-drag ratio that maximizes the horizontal distance traveled per unit of altitude lost during unpowered flight. Best glide speed is typically near the maximum L/D point on the polar curve and varies with aircraft weight and configuration. It is a critical parameter for emergency descent planning, forced landing procedures, and sailplane cross-country flying.}{catAD}

\dictentry{Blended Wing Body}{An aircraft configuration in which the wing and fuselage are smoothly integrated into a single lifting body, eliminating the conventional tube-and-wing distinction. The blended wing body offers significant reductions in wetted area and parasite drag, improving aerodynamic efficiency and fuel economy. It presents challenges in cabin pressurization, emergency egress, and airport compatibility.}{catAD}
\dictentry{Blockage}{The ratio of the frontal area of a wind tunnel model to the cross-sectional area of the test section, expressed as a percentage. High blockage ratios cause significant flow acceleration around the model and wall interference effects that distort the measured aerodynamic data. Typical guidelines require blockage below 1\% for closed-wall tunnels, with larger models requiring wall interference corrections.}{catAD}
\dictentry{Boundary Layer}{The thin region of flow immediately adjacent to a solid surface where viscous effects are dominant and the flow velocity increases from zero at the wall (no-slip condition) to the freestream value. The boundary layer can be laminar, transitional, or turbulent, and its state profoundly affects skin friction drag, heat transfer, and flow separation. Understanding boundary layer behavior is fundamental to aerodynamic design and performance prediction.}{catAD}
\dictentry{Boundary Layer Control}{Techniques used to modify the state, thickness, or behavior of the boundary layer to improve aerodynamic performance. Methods include suction through porous surfaces to maintain laminar flow, blowing to energize separated regions, and vortex generators to promote turbulent mixing. Boundary layer control is used to delay stall, reduce drag, and extend the operating envelope of wings and inlets.}{catAD}
\dictentry{Boundary Layer Separation}{The phenomenon in which the boundary layer detaches from the surface of a body when it encounters an adverse pressure gradient too strong for the near-wall flow to overcome. Separation results in a large wake of recirculating, low-pressure flow that dramatically increases pressure drag and can cause loss of lift and control effectiveness. It is the primary mechanism behind aerodynamic stall on wings.}{catAD}
\dictentry{Boundary Layer Transition}{The process by which flow in the boundary layer changes from an ordered laminar state to a chaotic turbulent state, triggered by Tollmien-Schlichting waves, surface roughness, freestream turbulence, or other disturbances. Transition typically increases skin friction drag by a factor of five to ten and significantly augments convective heat transfer. Predicting and controlling transition location is a major focus of laminar flow aircraft design.}{catAD}
\dictentry{Bow Shock}{A detached, curved shock wave that stands off from the leading edge or stagnation region of a blunt body traveling at supersonic or hypersonic speed. The bow shock decelerates the flow to subsonic speeds at the stagnation point and produces the highest pressures and temperatures in the flow field. The standoff distance depends on the body's bluntness and the freestream Mach number.}{catAD}
\dictentry{Buffeting}{Unsteady aerodynamic loading and vibration of an aircraft structure caused by turbulent wake, separated flow, or shock wave oscillations impinging on aft surfaces. Buffeting limits the usable lift coefficient, degrades ride quality, and can cause structural fatigue if sustained. It is commonly encountered at high angles of attack approaching stall and in the transonic regime due to shock-induced separation.}{catAD}
\dictentry{Bump}{A localized surface protuberance used for passive or active flow control, such as managing shock-boundary layer interaction on transonic wings. Bumps can weaken shock strength, delay separation, and redistribute surface pressure to improve aerodynamic efficiency. The concept has been applied to natural laminar flow wing designs and transonic airfoil optimization.}{catAD}
\dictentry{Buoyancy}{The upward force exerted by a fluid on a body immersed in it, equal to the weight of the fluid displaced by the body (Archimedes' principle). In aerodynamics, buoyancy is the principle behind lighter-than-air vehicles such as balloons and airships, and is proportional to the $\rho$ difference between the lifting gas and the surrounding air.}{catAD}
\dictentry{Camber}{The asymmetric curvature of an airfoil's mean line relative to the chord line, which causes the airfoil to generate lift at zero angle of attack. Positive camber increases the lift coefficient and shifts the zero-lift angle of attack to a negative value, improving low-speed performance. Excessive camber increases pitching moment and drag, so designers balance camber against trim and structural constraints.}{catAD}
\dictentry{Camber Line}{The locus of points equidistant from the upper and lower surfaces of an airfoil, defining the mean curvature of the section. The camber line shape determines the airfoil's lift-producing characteristics, including the zero-lift angle of attack and the pitching moment coefficient. It is a fundamental geometric parameter used in thin airfoil theory and airfoil design.}{catAD}
\dictentry{Canard}{A small auxiliary lifting surface mounted ahead of the main wing, used to provide pitch control, improve longitudinal stability, or enhance overall aerodynamic performance. Canards can generate positive lift in cruise, reducing trim drag compared to conventional tail-aft configurations, and are used on aircraft such as the Eurofighter Typhoon and Saab Gripen. The canard's proximity to the main wing creates complex aerodynamic interactions that must be carefully managed.}{catAD}
\dictentry{Carbon Footprint}{The total amount of carbon dioxide (CO2) emissions attributable to an aircraft or aviation operation over a given period or mission. Aviation carbon footprint reduction strategies include sustainable aviation fuels, improved aerodynamics, lightweight structures, and optimized flight operations.}{catAD}
\dictentry{Catalytic Heating}{The additional heating that occurs when atomic species (oxygen and nitrogen) recombine on a spacecraft surface during re-entry, releasing chemical energy. Catalytic heating depends on the surface material's recombination efficiency and can significantly increase total heat flux.}{catAD}
\dictentry{Center of Pressure}{The point on an airfoil or body at which the resultant aerodynamic force can be considered to act, such that there is no net pitching moment about that point. The center of pressure shifts aft with increasing angle of attack for cambered airfoils, which has important implications for longitudinal stability. It is distinct from the aerodynamic center, about which the moment remains constant.}{catAD}
\dictentry{Chevron Nozzle}{A jet engine exhaust nozzle with a serrated or sawtooth trailing edge that promotes mixing between the hot exhaust jet and the ambient air. Chevrons reduce jet noise by breaking up large-scale turbulent structures and shifting the noise spectrum to higher frequencies.}{catAD}
\dictentry{Chine}{A sharp longitudinal edge or strake-like feature on a fuselage or body that generates vortical lift at high angles of attack. Chines are commonly found on lifting bodies, stealth aircraft, and high-speed marine craft, where they contribute both to aerodynamic lift and to directional stability. At supersonic speeds, chines can also serve as wave riders that trap high-pressure air beneath the body.}{catAD}
\dictentry{Choked Flow}{A condition in a converging-diverging duct where the flow at the minimum cross-sectional area (throat) reaches sonic velocity (Mach 1), limiting the maximum mass flow rate through the duct regardless of further decreases in downstream pressure. Choked flow occurs when the pressure ratio across the duct exceeds a critical value. It is a key consideration in wind tunnel design, rocket nozzle operation, and gas turbine performance.}{catAD}
\dictentry{Chord}{The straight-line distance measured from the leading edge to the trailing edge of an airfoil section, used as a reference length for aerodynamic calculations. The chord varies along the span of a tapered wing and is used to define key parameters such as angle of attack, Reynolds number, and lift and drag coefficients. Mean aerodynamic chord (MAC) is the weighted average used for moment and stability calculations.}{catAD}
\dictentry{Chord Line}{The straight imaginary line connecting the leading edge stagnation point to the trailing edge of an airfoil, used as the reference datum for measuring angle of attack. The chord line is also the baseline for defining camber, thickness distribution, and the geometric properties of the airfoil section. It coincides with the chord for symmetric airfoils but differs for cambered sections.}{catAD}
\dictentry{Circulation}{The line integral of the velocity vector around a closed curve in a fluid, which by the Kutta-Joukowski theorem is directly proportional to the lift generated by a two-dimensional airfoil. Circulation quantifies the net rotational tendency of the flow around an airfoil and arises from the asymmetric flow speeds on upper and lower surfaces. It is a foundational concept in classical aerodynamic theory.}{catAD}
\dictentry{Clean Configuration}{The flight configuration of an aircraft in which all high-lift devices (flaps, slats) and landing gear are fully retracted, minimizing drag and maximizing airspeed for a given power setting. Clean configuration is used for cruise, climb, and high-speed flight.}{catAD}
\dictentry{Closed Circuit Tunnel}{A wind tunnel configuration in which the test airflow is recirculated through a continuous return loop, consisting of a test section, diffuser, turning vanes, settling chamber, contraction cone, and fan. Closed circuit tunnels offer better flow quality and energy efficiency than open circuit designs because the air is conditioned and reused. They are the standard configuration for most modern subsonic and transonic research facilities.}{catAD}
\dictentry{Cold Structure}{An aerospace structural design approach in which the primary load-bearing structure is protected from high temperatures by external insulation, maintaining it at relatively low temperatures. Cold structures allow the use of conventional materials like aluminum but add insulation weight.}{catAD}
\dictentry{Combustion Noise}{The noise generated by unsteady heat release and turbulent mixing in the combustion chamber of a gas turbine engine. Combustion noise contributes to the overall engine noise signature and is influenced by fuel-air mixing and combustor geometry.}{catAD}
\dictentry{Compressible Flow}{Fluid flow in which significant changes in gas $\rho$ occur in response to pressure variations, typically becoming important at Mach numbers above approximately 0.3. Compressibility effects include the formation of shock waves, expansion fans, and $\rho$ gradients that fundamentally alter the aerodynamic behavior of bodies. Compressible flow analysis requires solving the full continuity, momentum, and energy equations with variable $\rho$.}{catAD}
\dictentry{Compressor}{A rotating turbomachinery component that increases the pressure and temperature of a working fluid (typically air) by doing work on it through a series of rotating and stationary blade rows. Axial compressors are used in modern jet engines for their high mass flow rates and efficiency, while centrifugal compressors are simpler and more robust. Compressor performance is characterized by pressure ratio, efficiency, and surge margin.}{catAD}
\dictentry{Computational Fluid Dynamics}{A branch of fluid mechanics that uses numerical methods and algorithms to solve and analyze problems involving fluid flows, by discretizing the governing equations (Navier-Stokes, Euler, or simplified forms) over a computational grid. CFD enables detailed visualization of flow fields, force predictions, and parametric studies that would be prohibitively expensive through experimentation alone. It is now an indispensable tool in aerospace design, complementing wind tunnel testing.}{catAD}
\dictentry{Continuity Equation}{The mathematical statement of conservation of mass in fluid flow, requiring that the rate of mass entering a control volume equals the rate of mass leaving plus any accumulation within. For steady, incompressible flow it simplifies to the divergence of velocity being zero. The continuity equation is one of the fundamental governing equations used in all branches of fluid mechanics and aerodynamics.}{catAD}
\dictentry{Contra-Rotating Propeller}{A propulsion configuration consisting of two coaxial propellers mounted in tandem that rotate in opposite directions, recovering the swirl energy lost by the forward propeller. This arrangement increases propulsive efficiency and thrust for a given power input compared to a single propeller, particularly at high disk loading conditions. Contra-rotating propellers are used on some turboprop engines and marine propulsion systems.}{catAD}
\dictentry{Contraction Cone}{The converging section of a wind tunnel that accelerates the flow from the settling chamber to the test section. The contraction ratio (area ratio from inlet to outlet) determines the velocity increase and flow uniformity, with higher ratios producing lower turbulence levels.}{catAD}
\dictentry{Contrail}{A condensation trail composed of ice crystals that forms in the wake of an aircraft when hot, humid exhaust gases mix with the cold ambient air at high altitudes. Contrails can persist and spread into cirrus-like cloud sheets, contributing to radiative forcing and climate impact. Their formation depends on atmospheric temperature, humidity, and the water vapor output of the engines.}{catAD}
\dictentry{Control Reversal}{An aeroelastic phenomenon in which the deflection of a control surface produces an aerodynamic response opposite to the intended effect, typically occurring at high speeds due to wing twisting. As the control surface deflects, the resulting aerodynamic moment twists the wing in the opposite direction, overpowering the direct control effect. Control reversal sets an upper speed limit for conventional control systems and is mitigated through structural stiffness or predictive control laws.}{catAD}
\dictentry{Convective Heating}{The component of aerodynamic heating resulting from the transfer of thermal energy from the hot boundary layer gas to the vehicle surface through conduction and convection. Convective heating dominates at moderate re-entry velocities and is influenced by boundary layer state (laminar or turbulent).}{catAD}
\dictentry{Coordinated Turn}{A banked turn in which the aircraft's longitudinal axis is aligned with the relative wind, producing no sideslip, with lift and centripetal force properly balanced. In a coordinated turn, the ball in the turn coordinator remains centered, indicating efficient use of aerodynamic forces. Pilots achieve coordination through appropriate rudder input to match the bank angle and turn rate.}{catAD}
\dictentry{Corner Vanes}{Curved guide vanes installed in the corners of closed-circuit wind tunnels to turn the flow with minimal pressure loss and flow separation. They are essential for maintaining flow quality and reducing the power required to drive the tunnel.}{catAD}
\dictentry{Critical Angle of Attack}{The angle of attack at which an airfoil or wing reaches its maximum lift coefficient, beyond which lift decreases due to large-scale flow separation and stall. The critical angle is typically between 12 and 20 degrees depending on airfoil shape, Reynolds number, and leading edge geometry. It represents the upper limit of the usable lift range and is a key parameter in flight envelope definition.}{catAD}
\dictentry{Critical Mach Number}{The freestream Mach number at which the flow over some point on the airframe first reaches local sonic velocity (Mach 1), marking the onset of compressibility effects such as shock waves and drag rise. For conventional subsonic airfoils, the critical Mach number is typically around 0.7 to 0.85. Supercritical airfoils are designed to raise the critical Mach number by flattening the upper surface to delay shock formation.}{catAD}
\dictentry{Cruciform Tail}{A tail empennage configuration in which the horizontal and vertical stabilizers intersect at approximately 90 degrees, forming a cross-like planform when viewed from behind the aircraft. The cruciform arrangement separates the horizontal tail from the wing wake and engine exhaust while keeping structural weight moderate. It is commonly used on business jets, regional aircraft, and some military transports.}{catAD}
\dictentry{Cryogenic Wind Tunnel}{A wind tunnel that uses cryogenically cooled nitrogen gas (typically at temperatures around 100 K) as the test medium to achieve high Reynolds numbers at moderate model scales and tunnel speeds. Cooling the gas increases its $\rho$ and decreases its viscosity, allowing full-scale Reynolds number simulation in smaller facilities. The National Transonic Facility at NASA Langley is a prominent example.}{catAD}
\dictentry{Decalage}{The angular difference between the chord lines of the upper and lower wings of a biplane, measured in degrees. Positive decalage (upper wing at a higher angle of incidence) results in greater lift from the upper wing and is used to tailor stall characteristics and longitudinal stability.}{catAD}
\dictentry{Deceleron}{A split aileron in which the upper and lower halves can be deflected independently, allowing it to function simultaneously as a roll control surface and as a drag-inducing airbrake. When both halves deflect symmetrically, they act as a speed brake; when deflected differentially, they produce a rolling moment. Decelerons are used on some military and experimental aircraft for enhanced flight control versatility.}{catAD}
\dictentry{Delta Wing}{A wing with a triangular planform characterized by a highly swept leading edge and a straight or slightly swept trailing edge. Delta wings offer low wave drag at supersonic speeds and large internal volume for fuel and structure, making them popular for high-speed fighter aircraft such as the Concorde and the F-16 in its close-coupled canard-delta configuration. At low speeds they rely on vortex lift from leading-edge separation.}{catAD}
\dictentry{Detached Shock}{A curved shock wave that is not attached to the leading edge of a body but instead stands off at a finite distance, forming ahead of blunt bodies in supersonic or hypersonic flow. The detached shock decelerates the flow to subsonic speeds over a region upstream of the body and produces extreme temperatures and pressures in the stagnation region. The standoff distance increases with body bluntness and decreases with increasing Mach number.}{catAD}
\dictentry{Diffuser}{A duct section with an expanding cross-sectional area that decelerates a high-speed flow and converts kinetic energy into static pressure recovery. Diffusers are essential components of supersonic wind tunnels, jet engine inlets, and ramjet/scramjet systems, where they manage shock structures and deliver uniform subsonic flow to downstream components. Efficient pressure recovery in the diffuser directly impacts overall system performance.}{catAD}
\dictentry{Diffuser Section}{The expanding section of a wind tunnel located downstream of the test section, designed to decelerate the flow and recover static pressure. A well-designed diffuser minimizes pressure losses and improves overall tunnel efficiency.}{catAD}
\dictentry{Dihedral}{An upward angular inclination of the wing from root to tip relative to the lateral axis, which provides lateral (roll) stability by generating a restoring rolling moment when the aircraft sideslips. The greater the dihedral angle, the stronger the self-correcting tendency in roll after a disturbance. Low-wing aircraft typically require more dihedral than high-wing designs due to the different rolling moments produced by sideslip on the fuselage.}{catAD}
\dictentry{Dimple}{A small, localized surface indentation used to promote earlier boundary layer transition from laminar to turbulent, or to reduce drag by modifying the pressure distribution, as famously demonstrated on golf balls. On golf balls, dimples trip the boundary layer turbulent, which delays separation and dramatically reduces pressure drag, allowing the ball to travel farther. The concept has been explored for drag reduction on other bluff and streamlined bodies.}{catAD}
\dictentry{Disk Loading}{The average pressure change across an actuator disk such as a propeller or rotor, defined as the thrust divided by the disk area. Low disk loading is characteristic of rotors and high-lift propellers, while high disk loading is typical of high-speed propellers and ducted fans.}{catAD}
\dictentry{Displacement Thickness}{A measure of the outward displacement of the freestream streamline caused by the retardation of flow within the boundary layer due to viscous effects. It represents the distance by which a solid boundary would need to be displaced in an inviscid flow to produce the same mass flow deficit as the actual boundary layer. The displacement thickness is used in boundary layer analysis and wind tunnel wall interference corrections.}{catAD}
\dictentry{Divergence}{An aeroelastic instability in which a static aerodynamic load causes progressive structural deformation that increases the aerodynamic load in a feedback loop, leading to eventual structural failure. Divergence occurs when the aerodynamic center is ahead of the structural elastic axis and the dynamic pressure exceeds a critical value. It sets a speed limit for flexible wings and is prevented through sufficient torsional stiffness or sweep.}{catAD}
\dictentry{DNS}{Direct Numerical Simulation, a computational fluid dynamics approach that solves the full Navier-Stokes equations without any turbulence model, resolving all scales of motion from the largest energy-containing eddies down to the smallest Kolmogorov dissipative scales. DNS provides benchmark-quality data for turbulence research but is extremely computationally expensive, limiting it to low Reynolds numbers and simple geometries. It is used to validate turbulence models and study fundamental flow physics.}{catAD}
\dictentry{Dogtooth}{A small, sharp protuberance or notch added to the leading edge of a swept wing, designed to generate a streamwise vortex that energizes the boundary layer and delays spanwise flow separation at high angles of attack. The dogtooth vortex mimics the effect of a leading-edge extension (LEX) on a smaller scale. It was used on the Hawker Hunter and other early swept-wing fighters to improve high-speed handling.}{catAD}
\dictentry{Doublet}{A fundamental potential flow element formed by bringing a source and a sink of equal strength infinitesimally close together while increasing their strength, producing a flow pattern that decays as the inverse square of distance. The doublet is used as a building block in panel methods and conformal mapping to construct analytic solutions for flow around cylinders, airfoils, and other bodies. Combined with uniform flow, a doublet produces flow around a circular cylinder.}{catAD}
\dictentry{Downwash}{The induced downward velocity component imparted to the air by a lifting wing, which tilts the local lift vector backward and creates induced drag. Downwash is strongest near the wingtips where trailing vortices form and decreases with increasing distance behind the wing. The downwash distribution is directly related to the spanwise lift distribution and is a central element of lifting line theory.}{catAD}
\dictentry{Drag}{The aerodynamic force component acting parallel to and opposing the direction of motion of a body through the air. Drag consists of parasitic components (skin friction and pressure drag) and induced drag from lift production, and must be overcome by thrust for sustained flight. Minimizing drag is one of the primary objectives of aerodynamic design, as it directly affects fuel consumption, range, and speed.}{catAD}
\dictentry{Drag Coefficient}{A dimensionless parameter that quantifies the drag of a body relative to the dynamic pressure and reference area, defined as the drag force divided by the product of dynamic pressure and reference area. It encapsulates the effects of body shape, surface roughness, angle of attack, and Reynolds number into a single number for a given flow condition. Drag coefficient is the primary metric used to compare the aerodynamic efficiency of different shapes.}{catAD}
\dictentry{Dutch Roll}{A coupled lateral-directional oscillatory motion in which the aircraft simultaneously yaws and rolls with a characteristic frequency and damping ratio, resembling the motion of a skater weaving on ice. The oscillation is caused by the coupling between directional stability (weathercock effect) and lateral stability (dihedral effect). Excessive Dutch roll can be uncomfortable for passengers and is mitigated through yaw dampers and careful design of tail size and wing dihedral.}{catAD}
\dictentry{Dynamic Pressure}{The kinetic energy per unit volume of a fluid flow, defined as q = 0.5 * $\rho V^2$. Dynamic pressure represents the pressure rise when flow is brought to rest isentropically from its current velocity to stagnation, and is the fundamental scaling parameter for aerodynamic forces and moments. Aerodynamic force coefficients are defined relative to dynamic pressure and reference area.}{catAD}

\dictentry{Dynamic Stability}{The long-term behavior of an aircraft's motion following a disturbance, describing whether oscillations or displacements grow, decay, or remain constant over time. A dynamically stable aircraft will see its transient motions diminish after the initial static stability response. Dynamic stability analysis involves solving the equations of motion for the aircraft's short period, phugoid, Dutch roll, and spiral modes.}{catAD}
\dictentry{Dynamic Stall}{A transient aerodynamic phenomenon occurring when an airfoil undergoes rapid changes in angle of attack, causing a delayed stall with the formation of a leading-edge vortex that temporarily augments lift beyond the static stall limit. The vortex eventually sheds and convects downstream, producing a sharp drop in lift and a large nose-down pitching moment. Dynamic stall is critical for helicopter rotor blade performance and maneuvering fighter aircraft.}{catAD}
\dictentry{Dynamic Testing}{Aerodynamic or structural testing conducted under time-varying conditions, including oscillatory model motion, pulsating flow, or transient loads, to measure unsteady aerodynamic responses and frequency-dependent behavior. Dynamic testing captures phenomena such as flutter derivatives, indicial responses, and aerodynamic admittance that cannot be obtained from steady measurements. It is essential for validating aeroelastic predictions and flight dynamics models.}{catAD}
\dictentry{Eddy}{A local swirling motion of fluid within a turbulent flow, characterized by vorticity and a distinct velocity field that differs from the mean flow. Eddies span a wide range of scales from the large energy-containing eddies down to the Kolmogorov microscales where viscous dissipation occurs. The energy cascade transfers kinetic energy from large to small eddies until it is dissipated as heat.}{catAD}

\dictentry{Electric Propulsion}{Aircraft propulsion systems that use electric motors powered by batteries, fuel cells, or hybrid generators to drive propellers or fans. Electric propulsion enables distributed propulsion configurations, reduces local emissions, and opens new vehicle architectures for urban air mobility.}{catAD}
\dictentry{Elevator}{A hinged control surface mounted on the horizontal tail of a conventional aircraft, used to control pitch attitude and thereby angle of attack and airspeed. Downward deflection of the elevator increases tail lift, pushing the nose down, while upward deflection decreases tail lift, raising the nose. The elevator is one of the three primary flight controls and works in conjunction with trim systems to maintain steady flight.}{catAD}
\dictentry{Elevon}{A combined elevator and aileron control surface mounted on the trailing edge of a tailless or flying-wing aircraft, providing both pitch and roll control from a single surface. Elevons deflect differentially for roll control and symmetrically for pitch control, eliminating the need for separate horizontal tail and aileron surfaces. This configuration reduces weight, drag, and complexity in tailless designs such as the B-2 Spirit and Concorde.}{catAD}
\dictentry{Elliptical Wing}{A wing planform in which the spanwise lift distribution is elliptical, producing the minimum possible induced drag for a given total lift and wingspan according to Prandtl's lifting line theory. The elliptical distribution ensures uniform downwash across the span, eliminating induced drag due to non-optimal lift distribution. While aerodynamically ideal, elliptical wings are structurally complex and expensive to manufacture, so practical designs approximate the elliptical distribution with taper and twist.}{catAD}
\dictentry{Emission Index}{A quantitative measure of pollutant mass produced per unit mass of fuel burned in an aircraft engine, expressed in grams of pollutant per kilogram of fuel. Emission indices for NOx, CO, unburned hydrocarbons, and soot are used to assess engine environmental performance.}{catAD}
\dictentry{Endplate}{A flat vertical plate attached at the wingtip of an aircraft or hydrofoil, designed to reduce induced drag by acting as a barrier that prevents high-pressure air from the lower surface from spilling around the tip to the low-pressure upper surface. Endplates effectively increase the aspect ratio without increasing the physical wingspan, which is valuable for aircraft with span constraints. They are commonly used on hydrofoils, racing cars, and some UAV designs.}{catAD}
\dictentry{Endurance}{The maximum time an aircraft can remain airborne on a given fuel or energy supply, achieved by minimizing the fuel flow rate per unit time. Endurance is optimized by flying at the speed for minimum power required (maximum L/D for propeller aircraft, different for jets). Maximum endurance speed is typically lower than maximum range speed and is used for loiter, surveillance, and holding patterns.}{catAD}

\dictentry{Energy Equation}{The mathematical statement of conservation of energy applied to fluid flow, relating changes in internal energy, kinetic energy, and potential energy to work done by pressure forces and heat transfer. For adiabatic, inviscid flow it reduces to the total enthalpy remaining constant along a streamline. The energy equation is essential for analyzing compressible flows, shock waves, and stagnation properties.}{catAD}
\dictentry{Enstrophy}{A quantity in fluid dynamics that measures the integrated square of the vorticity magnitude, representing the dissipation of kinetic energy in turbulent flows. Enstrophy is particularly useful in the study of turbulence and is directly related to the turbulent kinetic energy dissipation rate.}{catAD}
\dictentry{Euler Equations}{The governing equations of inviscid fluid dynamics, representing the conservation of mass, momentum, and energy for a frictionless fluid, obtained by neglecting the viscous stress terms in the Navier-Stokes equations. Euler equations are applicable to high Reynolds number flows away from boundary layers and are used to model shock waves, expansion fans, and external aerodynamics at supersonic speeds. They form the basis for many computational methods in compressible flow analysis.}{catAD}
\dictentry{Euler Number}{A dimensionless parameter defined as the ratio of pressure forces to inertial forces, Eu = $\Delta p$ / ($\rho V^2$). It characterizes the relationship between pressure drop and dynamic pressure in fluid systems and is used in cavitation analysis, pressure loss calculations, and compressible flow similarity. For incompressible flow, the Euler number is closely related to the pressure coefficient.}{catAD}

\dictentry{Expansion Fan}{A continuous region of expanding flow that forms when a supersonic flow turns around a convex corner or expands over a convex surface, with each infinitesimal turning angle generating a weak Mach wave. Across the expansion fan, the Mach number increases, while pressure, temperature, and $\rho$ decrease isentropically. Expansion fans are the supersonic counterpart to compression shocks and are fundamental to Method of Characteristics solutions.}{catAD}
\dictentry{Expansion Wave}{A weak, isentropic wave in supersonic flow that slightly expands and accelerates the flow while reducing pressure and temperature, generated when the flow encounters a convex corner or surface. Unlike shock waves, expansion waves are continuous and reversible, with no entropy increase. A series of expansion waves emanating from a sharp convex corner forms an expansion fan described by the Prandtl-Meyer function.}{catAD}
\dictentry{External Balance}{A force measurement balance mounted outside the wind tunnel test section, connected to the model through slender struts or wires that pass through the tunnel wall. External balances isolate the sensitive measuring elements from the harsh test section environment (temperature, vibration, pressure), improving accuracy and ease of calibration. However, the support struts introduce additional interference that must be accounted for in data reduction.}{catAD}
\dictentry{Fan Noise}{The noise produced by the fan stage of turbofan engines, comprising tonal noise at the blade passage frequency and broadband noise from turbulent interactions. Fan noise is a significant contributor to aircraft noise during takeoff and approach.}{catAD}
\dictentry{Fence}{A thin, vertical surface mounted on the upper or lower surface of a wing to control or block spanwise flow, preventing the migration of boundary layer air toward the wingtip. Fences improve aileron effectiveness at high angles of attack and prevent asymmetric stall by maintaining attached flow over the outboard wing sections. They are commonly found on swept-wing aircraft where spanwise flow is a significant concern.}{catAD}
\dictentry{Finite Element Method}{A numerical technique for finding approximate solutions to boundary value problems by subdividing a complex domain into smaller, simpler elements (typically triangles or tetrahedra) over which simple polynomial functions approximate the solution. In aerodynamics and structures, FEM is widely used for structural analysis, flutter prediction, and coupled fluid-structure interaction simulations. Its strength lies in handling complex geometries and heterogeneous materials.}{catAD}
\dictentry{Finite Volume Method}{A numerical method for solving partial differential equations by discretizing the domain into control volumes and applying conservation laws in integral form to each volume, ensuring that fluxes are balanced across cell interfaces. The finite volume method naturally conserves quantities like mass, momentum, and energy, making it the preferred approach for computational fluid dynamics codes. It handles complex geometries and unstructured meshes effectively.}{catAD}
\dictentry{Fish-Tail Shock}{A distinctive shock wave pattern that resembles a fish tail when observed in schlieren or shadowgraph images of transonic or supersonic flow over nozzles, airfoils, or blunt bodies. The pattern results from the interaction of oblique shocks and slip lines emanating from a trailing edge or nozzle lip. Fish-tail shocks are associated with wave drag and are studied in nozzle design and transonic airfoil optimization.}{catAD}
\dictentry{Flap}{A hinged or sliding high-lift device mounted on the trailing edge of a wing that increases the effective camber and/or wing area when deployed, boosting the maximum lift coefficient for takeoff and landing. Common types include plain flaps, split flaps, slotted flaps, and Fowler flaps, each offering different trade-offs between lift increment, drag, and mechanical complexity. Flaps are retracted during cruise to restore the clean wing aerodynamic shape.}{catAD}
\dictentry{Flapperon}{A dual-function control surface that operates as a conventional flap for increasing lift during takeoff and landing, and as an aileron for roll control during cruise. The flapperon combines the trailing edge flap and aileron into a single continuous surface, reducing the number of control surfaces and simplifying the wing structure. It is used on some light aircraft and UAVs where weight and complexity reduction are priorities.}{catAD}
\dictentry{Flexible TPS}{A lightweight, conformable thermal protection material made of ceramic or polymer fibers, used for areas with moderate heating on re-entry vehicles. Flexible TPS blankets can be draped over curved surfaces and are easier to install than rigid tiles.}{catAD}
\dictentry{Flow Angularity}{The deviation of the local flow direction from the nominal tunnel axis or freestream direction within a wind tunnel test section, caused by model wake effects, wall interference, or imperfections in the flow conditioning system. Flow angularity introduces errors in angle of attack and sideslip measurements and must be minimized through careful tunnel design and calibration. Typical acceptable limits are on the order of 0.1 to 0.2 degrees.}{catAD}
\dictentry{Flow Attachment}{The condition in which a boundary layer remains in continuous contact with a solid surface, following its contour without separation. Attached flow produces predictable pressure distributions, low pressure drag, and effective lift generation. Maintaining flow attachment at high angles of attack is a primary goal of high-lift system design and flow control methods.}{catAD}

\dictentry{Flow Meter}{An instrument for measuring the mass flow rate or volumetric flow rate of a fluid in a pipe or duct. In aerodynamic testing, flow meters measure air supply rates for wind tunnels, engine test cells, and boundary layer suction/blowing systems. Types include orifice plates, venturi meters, thermal mass flow meters, and ultrasonic flow meters.}{catAD}

\dictentry{Flow Separation}{The detachment of the boundary layer from a solid surface when it encounters an adverse pressure gradient too strong for the near-wall momentum to overcome, resulting in a region of reversed or recirculating flow downstream of the separation point. Separation causes a large increase in pressure drag, loss of lift, and unsteady wake formation. It is the primary mechanism behind aerodynamic stall and limits the maximum usable angle of attack of wings.}{catAD}
\dictentry{Flutter}{A dangerous aeroelastic instability in which aerodynamic forces, structural elasticity, and inertial effects couple to produce self-excited oscillations that can grow exponentially and lead to catastrophic structural failure within seconds. Flutter occurs when the energy extracted from the airflow exceeds the energy dissipated by structural damping, and its onset speed defines a critical flight envelope limit. Prevention requires careful aeroelastic analysis, ground vibration testing, and flight testing to ensure adequate flutter margins.}{catAD}
\dictentry{Flying Wing}{An aircraft configuration in which the entire airframe serves as a lifting surface, with no distinct fuselage or tail, eliminating the parasitic drag and weight of non-lifting components. Flying wings offer superior aerodynamic efficiency and low radar cross-section, making them attractive for long-range bombers, transports, and stealth aircraft such as the B-2 Spirit. The configuration presents challenges in longitudinal stability, pitch control, and volume for payload and systems.}{catAD}
\dictentry{Force Balance}{A multi-component measuring device installed inside or connected to a wind tunnel model that simultaneously resolves the total aerodynamic force into its lift, drag, and side force components, as well as pitching, rolling, and yawing moments. Force balances use arrays of strain gauges arranged in specific bridge circuits to decouple the six force and moment components. Accuracy, stiffness, and thermal stability are critical design requirements for reliable wind tunnel measurements.}{catAD}
\dictentry{Form Drag}{The component of aerodynamic drag arising from the pressure distribution around a body, caused by the difference between the high pressure on the forward-facing surfaces and the low pressure in the separated wake behind the body. Form drag depends strongly on the body's shape and frontal area, with streamlined bodies having much lower form drag than blunt bodies. Reducing form drag involves smoothing the body contour, tapering the aft section, and delaying flow separation.}{catAD}
\dictentry{Forward Swept Wing}{A wing configuration in which the leading edge is angled forward from root to tip, rather than rearward as in conventional swept-wing designs. Forward sweep delays the onset of tip stall and maintains aileron effectiveness at high angles of attack because the boundary layer flows inward toward the root rather than outward toward the tips. However, forward sweep introduces aeroelastic divergence problems that require very stiff wing structures or active load alleviation.}{catAD}
\dictentry{Friction Drag}{The component of aerodynamic drag caused by viscous shear stresses acting tangentially on the surface of a body as the boundary layer flows over it. Friction drag depends on the wetted area, boundary layer state (laminar or turbulent), and Reynolds number, with turbulent boundary layers producing roughly five times more friction drag than laminar ones. Laminar flow control techniques aim to maintain laminar boundary layers over larger portions of the surface to reduce friction drag.}{catAD}
\dictentry{Froude Number}{A dimensionless parameter defined as the ratio of inertial forces to gravitational forces, $Fr = V / \sqrt{gL}$. It is the primary similarity parameter for free-surface flows, ship hydrodynamics, and flows with significant gravity effects. In aerodynamics, it is relevant for ground effect vehicles, seaplane operations, and atmospheric flows with buoyancy.}{catAD}

\dictentry{Full Model Testing}{Wind tunnel testing of a complete aircraft model including all major components (fuselage, wings, tail, engine nacelles, landing gear) at the appropriate geometric scale, capturing the full aerodynamic interactions between components. Full model tests provide the most realistic representation of the complete aircraft's aerodynamic characteristics, including interference effects between components. They require larger test sections and higher tunnel power than isolated component tests.}{catAD}
\dictentry{Geometric Twist}{The physical twisting of a wing structure from root to tip that changes the local angle of incidence along the span, typically implemented as washout (decreasing angle toward the tip). Geometric twist is used to tailor the spanwise lift distribution, delay tip stall, and improve stall characteristics. It is a permanent structural feature that must be designed and manufactured into the wing during construction.}{catAD}
\dictentry{Grid}{The discretized spatial mesh of cells, elements, or nodes that divides a continuous fluid domain into finite regions for numerical computation in CFD. Grid quality, including cell size, aspect ratio, skewness, and resolution in regions of high gradients, directly affects the accuracy and convergence of the numerical solution. Grids can be structured (regular topology) or unstructured (flexible topology), with hybrid approaches often used for complex geometries.}{catAD}
\dictentry{Grid Turbulence}{A approximately isotropic, homogeneous turbulent flow generated by passing a uniform mean flow through a grid of bars or perforated plates, which breaks up the mean flow into small-scale eddies. Grid turbulence decays with distance downstream of the grid and is widely used as a reference flow for turbulence research, turbulence model validation, and wind tunnel flow quality characterization. It provides a well-defined turbulence environment free from mean shear effects.}{catAD}
\dictentry{Ground Effect}{The phenomenon in which the induced drag of a wing is significantly reduced when the aircraft flies at a height above the surface comparable to or less than its wingspan, due to the ground interfering with the downwash and tip vortex system. Ground effect causes an increase in lift-to-drag ratio and makes the aircraft feel as though it is floating, requiring a different flare technique during landing. It is most pronounced for high-aspect-ratio wings and is exploited by ekranoplanes (ground-effect vehicles).}{catAD}
\dictentry{Gust Load}{A transient aerodynamic load imposed on an aircraft structure when it encounters an atmospheric gust, causing a rapid change in angle of attack and consequent change in lift. Gust loads are a critical design driver for aircraft structural sizing, particularly for wing bending moments and fatigue life, and are specified by certification regulations as discrete and continuous turbulence requirements. The response depends on aircraft speed, weight, wing loading, and flexibility.}{catAD}
\dictentry{H-Tail}{A tail empennage configuration in which two vertical fins are mounted at the ends of a single horizontal stabilizer, resembling the letter H when viewed from behind. The H-tail provides redundancy in directional control and is favored on multi-engine aircraft where an engine failure creates significant asymmetric thrust. It also places the vertical surfaces outside the wing wake and jet exhaust for improved effectiveness.}{catAD}
\dictentry{Half Model Testing}{A wind tunnel testing technique in which a half-span model of an aircraft is mounted flush against a splitter plate or tunnel wall, exploiting the symmetry plane to simulate a full model at double the Reynolds number for the same tunnel speed. Half model testing allows larger model scales and higher Reynolds numbers within a given test section size, improving measurement accuracy. The splitter plate must be carefully designed to minimize boundary layer and interference effects.}{catAD}
\dictentry{Heat Flux}{The rate of thermal energy transfer per unit area, expressed in watts per square meter (W/m2), at a surface exposed to aerodynamic heating. Heat flux is a critical parameter in the design of thermal protection systems and is composed of convective, radiative, and catalytic components.}{catAD}

\dictentry{High Altitude Flight}{Flight operations at altitudes above 50,000 feet, where low air $\rho$ requires high aerodynamic efficiency and specialized propulsion systems. High altitude flight includes stratospheric platforms for communications, surveillance, and atmospheric science.}{catAD}
\dictentry{High Speed Tunnel}{A wind tunnel designed for testing at transonic (Mach 0.8 to 1.2) or supersonic (Mach 1.2 to 5) speeds, requiring specialized features such as perforated or slotted test section walls to relieve wave reflection, and converging-diverging nozzle sections to accelerate the flow. High speed tunnels must handle the unique challenges of shock wave-boundary layer interactions and compressibility effects. They are essential for determining drag rise, shock location, and buffet boundaries.}{catAD}
\dictentry{Hoerner Tip}{A rounded, drooped wingtip shape designed by Dr. S.F. Hoerner to reduce induced drag by weakening the tip vortex and reducing the pressure difference between the upper and lower surfaces at the tip. The Hoerner tip is simpler and lighter than winglets while still providing measurable drag reduction, and is widely used on light aircraft, sailplanes, and some military aircraft. It represents an effective compromise between drag reduction and structural simplicity.}{catAD}
\dictentry{Honeycomb Straightener}{A flow conditioning device consisting of a bundle of small-diameter tubes or cells placed upstream of the contraction cone in a wind tunnel. It reduces turbulence and swirl by breaking up large-scale eddies and aligning the flow with the tunnel axis, improving flow quality in the test section.}{catAD}
\dictentry{Horseshoe Vortex}{The vortex system formed when a boundary layer approaching a surface-mounted obstacle (such as a wing-body junction or vertical fin) wraps around the obstacle's leading edge, creating a pair of counter-rotating vortices that extend downstream along the sides of the junction. The horseshoe vortex produces locally high shear stresses and heat transfer at the junction and can generate interference drag. It is a common feature in turbomachinery blade-endwall junctions and aircraft wing-root regions.}{catAD}
\dictentry{Hot Structure}{An aerospace structural design philosophy in which the primary load-bearing structure is allowed to reach high temperatures during flight, eliminating the need for heavy insulation. Hot structures are typically made of high-temperature alloys, titanium, or ceramic matrix composites.}{catAD}
\dictentry{Hot Wire Anemometry}{A flow measurement technique that uses a very thin electrically heated wire (typically tungsten or platinum, 1-10 micrometers in diameter) whose temperature, and hence electrical resistance, varies with the convective cooling effect of the surrounding airflow. By measuring the resistance change, the local flow velocity and its fluctuations can be determined with very high frequency response. Hot wire anemometry is widely used for turbulence measurements and boundary layer research.}{catAD}
\dictentry{Hush Kit}{A set of modifications applied to an aircraft engine or nacelle to reduce noise emissions, typically consisting of acoustic liners, exhaust mixers, and nozzle treatments. Hush kits allow older aircraft to meet modern noise certification standards.}{catAD}
\dictentry{Hybrid Propulsion}{A propulsion architecture combining a conventional thermal engine (turbine or piston) with an electric motor and battery system. Hybrid propulsion allows the thermal engine to operate at optimum efficiency while the electric motor provides peak power for takeoff and climb.}{catAD}
\dictentry{Hydrogen Propulsion}{The use of hydrogen as a fuel for aircraft propulsion, either burned in modified gas turbine engines or used in fuel cells to power electric motors. Hydrogen offers zero CO2 emissions at the point of use but presents challenges in storage volume, cryogenic handling, and infrastructure.}{catAD}
\dictentry{Hypersonic Flow}{Flow at Mach numbers greater than 5, where strong shock waves, extreme compression heating, viscous interaction, and real gas effects (dissociation, ionization) dominate the aerodynamic behavior. Hypersonic flight produces intense aerodynamic heating that requires specialized thermal protection systems and materials capable of withstanding extreme temperatures. The regime is relevant for ICBM re-entry, scramjet-powered vehicles, and planetary entry probes.}{catAD}
\dictentry{Hypersonic Wind Tunnel}{A wind tunnel capable of producing test flows at Mach numbers above 5, using high-pressure heated or high-enthalpy air to simulate the stagnation conditions of hypersonic flight. Hypersonic tunnels include blow-down facilities with heated driver gas, shock tunnels, and expansion tubes, each suited to different test durations and enthalpy ranges. They are used for aerodynamic force testing, heat transfer measurements, and scramjet inlet research.}{catAD}
\dictentry{Incompressible Flow}{A flow approximation in which the fluid $\rho$ is assumed constant regardless of pressure or velocity changes, valid for gas flows below approximately Mach 0.3 where compressibility effects are negligible. Incompressible flow analysis greatly simplifies the governing equations, allowing the use of potential flow theory, Bernoulli's equation, and simpler computational methods. Most low-speed aerodynamics, including general aviation and UAV flight, can be treated as incompressible.}{catAD}
\dictentry{Indraft Wind Tunnel}{An open-circuit wind tunnel configuration in which the fan is located downstream of the test section and draws ambient room air through the inlet, test section, and diffuser. Indraft tunnels are simpler and cheaper to build than closed-circuit designs, and the absence of flow conditioning upstream of the test section can reduce thermal contamination. However, they are sensitive to room disturbances and have higher power consumption for equivalent flow quality.}{catAD}
\dictentry{Induced Drag}{The component of drag that arises as an unavoidable consequence of generating lift on a finite wing, caused by the downwash from trailing vortices tilting the local lift vector backward. Induced drag is inversely proportional to the square of the airspeed and directly proportional to the square of the lift coefficient, making it dominant at low speeds and high angles of attack. Increasing wingspan (aspect ratio) is the most effective way to reduce induced drag.}{catAD}
\dictentry{Induced Velocity}{The local velocity perturbation induced at any point in the flow field by the bound and trailing vortex system of a lifting wing, most commonly referring to the downward velocity component (downwash) at the wing itself. Induced velocity varies along the span and is directly related to the local lift distribution through the Biot-Savart law. It is a key output of vortex lattice and lifting line methods and determines the induced angle of attack.}{catAD}
\dictentry{Inflatable TPS}{A deployable thermal protection system consisting of a flexible, inflatable aeroshell that decelerates a payload during atmospheric entry. Inflatable TPS enables larger drag areas than rigid aeroshells of the same mass, reducing peak heating and deceleration.}{catAD}
\dictentry{Infrared Thermography}{A non-contact, full-field temperature measurement technique that uses an infrared camera to detect thermal radiation emitted by a surface and convert it to a temperature map. In aerodynamic testing, infrared thermography is used to identify transition locations on wind tunnel models (by detecting the temperature difference between laminar and turbulent boundary layers), measure heat transfer rates, and detect structural hot spots. It provides excellent spatial resolution without disturbing the flow.}{catAD}

\dictentry{Interference Drag}{The additional drag component arising from the aerodynamic interaction between two or more components of an aircraft (such as wing-fuselage, pylon-nacelle, or strut-wing junctions) that exceeds the sum of the individual component drags measured in isolation. Interference effects can cause localized flow acceleration, shock formation, and early separation at junctions. Streamlining junction fairings and careful component spacing are primary methods for minimizing interference drag.}{catAD}
\dictentry{Internal Balance}{A force measurement balance mounted inside the hollow cavity of a wind tunnel model, with the sensing elements connected to the model's internal structure and the external support sting. Internal balances are protected from the flow environment and do not introduce external support interference, making them preferred for high-speed testing. They require careful installation, alignment, and thermal compensation to achieve accurate multi-component force measurements.}{catAD}
\dictentry{Inviscid Flow}{An idealized flow in which viscous effects (friction and thermal conductivity) are entirely neglected, allowing the fluid to slip freely along solid surfaces without forming a boundary layer. Inviscid flow theory, governed by the Euler equations, provides excellent predictions for external aerodynamics away from surfaces and is the basis for potential flow methods, lift prediction, and shock-expansion theory. It cannot predict skin friction, flow separation, or wake formation.}{catAD}
\dictentry{Irrotational Flow}{A flow field in which the local rotation (curl of the velocity vector) is zero everywhere, meaning that fluid particles translate and deform but do not rotate about their own axes. Irrotational flow is a key assumption in potential flow theory, enabling the velocity field to be expressed as the gradient of a scalar velocity potential that satisfies Laplace's equation. Most external aerodynamic flows outside the boundary layer and wake are approximately irrotational.}{catAD}
\dictentry{Jet Noise}{The noise generated by the high-velocity exhaust of a jet engine, resulting from turbulent mixing of the jet plume with the ambient air. Jet noise is a major component of aircraft noise during takeoff and is reduced through chevron nozzles and serrated ejectors.}{catAD}
\dictentry{Joukowski Airfoil}{An airfoil shape generated by applying the Joukowski conformal transformation to a circle, producing a family of shapes ranging from flat plates to cambered, rounded trailing-edge sections. The Joukowski airfoil was one of the first shapes for which an exact analytical solution of the inviscid flow field could be obtained, providing the theoretical basis for understanding lift generation and the Kutta condition. It remains a valuable pedagogical tool in aerodynamics education.}{catAD}
\dictentry{Keel Effect}{The aerodynamic phenomenon where lateral area located behind the center of gravity provides directional stability, analogous to the keel of a ship. In aircraft design, the keel effect is provided by the vertical tail and aft fuselage and contributes to weathercock stability.}{catAD}
\dictentry{Kutta Condition}{The physical condition applied in potential flow theory requiring that the flow leaves the sharp trailing edge of an airfoil smoothly, determining the unique value of circulation and hence the lift generated. Without the Kutta condition, inviscid theory would predict zero lift for a airfoil with a sharp trailing edge, which contradicts experiment. The condition implicitly accounts for the viscous effects in the boundary layer that enforce smooth trailing-edge flow.}{catAD}
\dictentry{Kutta-Joukowski Theorem}{The fundamental theorem of two-dimensional aerodynamics stating that the lift per unit span on a body in inviscid flow is equal to the product of the fluid $\rho$, the freestream velocity, and the circulation around the body. This elegant result, derived independently by Kutta and Joukowski in the early 1900s, provides the theoretical foundation for understanding how airfoils generate lift and relates the abstract concept of circulation to a measurable aerodynamic force.}{catAD}
\dictentry{Lambda Shock}{A distinctive lambda-shaped (lambda) shock wave structure observed in transonic flow over airfoils, consisting of a nearly normal shock segment on the upper surface that bifurcates into an oblique shock branch ahead and a slip line behind. The lambda shock forms when the normal shock interacts with the decelerating boundary layer, causing local separation. It influences pressure distribution, buffet onset, and drag characteristics in the transonic regime.}{catAD}
\dictentry{Laminar Flow}{A state of fluid motion characterized by smooth, orderly, parallel layers (laminae) of fluid with minimal mixing between adjacent layers, in which viscous forces dominate over inertial forces. Laminar boundary layers produce significantly less skin friction drag than turbulent ones but are susceptible to transition triggered by surface roughness, pressure gradients, or freestream disturbances. Maintaining laminar flow over larger portions of aircraft surfaces is a major goal of drag reduction research.}{catAD}
\dictentry{Laminar Flow Control}{Active or passive techniques designed to maintain laminar boundary layer flow over a greater portion of an aircraft's surface, thereby reducing skin friction drag. Active methods include suction through porous or perforated surfaces to remove the retarded boundary layer fluid, while passive methods involve careful surface shaping (natural laminar flow airfoils) and surface finish control. Laminar flow control has the potential to reduce total aircraft drag by 10-25 percent.}{catAD}
\dictentry{Landing Distance}{The horizontal distance required for an aircraft to decelerate from a defined height above the runway (typically 50 feet) to a full stop, including the air distance and ground roll. Landing distance depends on approach speed, weight, configuration, braking effectiveness, runway surface, wind, and pilot technique. Certification regulations require demonstrated landing distances with safety margins.}{catAD}

\dictentry{Laser Doppler Velocimetry}{An optical, non-intrusive flow measurement technique that determines local fluid velocity by detecting the Doppler frequency shift of laser light scattered by small seed particles carried in the flow. Two intersecting laser beams create an interference fringe pattern, and the frequency of the scattered light signal is directly proportional to the particle velocity. LDV provides high spatial and temporal resolution without disturbing the flow field.}{catAD}
\dictentry{Lateral Stability}{The tendency of an aircraft to resist and recover from uncommanded roll disturbances, primarily provided by wing dihedral, swept wings, and high-mounted fuselage lateral area. Lateral stability is assessed through the roll subsidence and Dutch roll modes and is characterized by the rolling moment due to sideslip derivative. Proper lateral stability is essential for acceptable handling qualities and coordinated flight.}{catAD}
\dictentry{LDV}{Laser Doppler Velocimetry (abbreviation). See Laser Doppler Velocimetry.}{catAD}

\dictentry{Leading Edge}{The forward-most edge of an airfoil, wing, fin, or other aerodynamic surface, which first encounters the oncoming airflow. The leading edge geometry (sharp, rounded, or drooped) has a profound influence on the airfoil's stall characteristics, maximum lift coefficient, and sensitivity to angle of attack and Mach number. Leading edge contamination by ice, insects, or surface damage can severely degrade aerodynamic performance.}{catAD}
\dictentry{Leading Edge Extension}{A strake or extension projecting forward and outboard from the wing root leading edge, designed to generate powerful streamwise vortices at high angles of attack that energize the flow over the outboard wing and delay stall. LEX devices are a key feature of modern fighter aircraft such as the F/A-18 and F-16, where they extend the usable angle of attack range well beyond the clean wing stall. The vortex system from the LEX also increases lift at high angles of attack.}{catAD}
\dictentry{Leading Edge TPS}{The thermal protection system applied to the leading edges of wings, fins, and control surfaces of re-entry vehicles and hypersonic aircraft. Leading edge TPS must withstand high heat fluxes while maintaining aerodynamic shape and structural integrity.}{catAD}
\dictentry{LES}{Large Eddy Simulation, a computational turbulence modeling approach in which the large energy-containing eddies are resolved directly on the computational grid while the smaller, more universal sub-grid scales are modeled using a sub-grid scale model. LES provides better accuracy than RANS for unsteady, separated, and vortical flows while being far less expensive than DNS. It is increasingly used for aerodynamic noise prediction, combustion modeling, and complex separated flows.}{catAD}
\dictentry{LEX Fence}{A small vertical fence or strake mounted on the upper surface of a leading edge extension to control the strength, position, and bursting characteristics of the vortex generated by the LEX. The fence prevents the vortex from migrating spanwise and bursting prematurely over the vertical tail, which could cause loss of directional control at high angles of attack. It is a key flow control feature on aircraft like the F/A-18 Hornet.}{catAD}
\dictentry{Lift}{The aerodynamic force component acting perpendicular to the relative wind (freestream direction) that opposes the weight of the aircraft and enables flight. Lift is generated by the pressure difference between the upper and lower surfaces of a wing, with lower pressure on top and higher pressure on the bottom, as described by circulation theory. The magnitude of lift depends on airspeed squared, air $\rho$, wing area, and lift coefficient.}{catAD}
\dictentry{Lift Coefficient}{A dimensionless parameter that quantifies the lift generated by an airfoil or wing relative to the dynamic pressure and reference area, defined as the lift force divided by the product of dynamic pressure and planform area. It varies with angle of attack, Mach number, Reynolds number, and wing geometry, and is the primary measure of lifting performance. The maximum lift coefficient determines the stall speed and low-speed flight characteristics.}{catAD}
\dictentry{Lift Distribution}{The spanwise variation of lift along a wing, determined by the wing planform shape, twist, camber distribution, and the influence of fuselage and tail surfaces. The lift distribution directly determines the induced drag, with an elliptical distribution producing the minimum induced drag for a given total lift and span. Designers use taper, sweep, and washout to approximate the elliptical distribution while meeting structural and manufacturing constraints.}{catAD}
\dictentry{Lifting Line Theory}{Prandtl's classical aerodynamic theory for finite wings that models the wing as a bound vortex line with trailing vortices forming a vortex sheet. The theory predicts spanwise lift distribution, induced drag, and downwash, and shows that an elliptical lift distribution produces minimum induced drag. It provides the theoretical basis for understanding finite wing aerodynamics and wing design.}{catAD}
\dictentry{Lifting Re-entry}{A re-entry trajectory that uses aerodynamic lift generated by the vehicle body or wings to control the descent path and reduce peak deceleration. Lifting re-entry, used by the Space Shuttle and proposed for crewed Mars missions, enables cross-range maneuvering and gentler deceleration.}{catAD}
\dictentry{Lifting Surface Theory}{An extension of lifting line theory that accounts for the finite chord of a wing by distributing singularities (such as vortex panels or doublets) over the entire wing surface rather than along a single line. Lifting surface methods can predict the aerodynamic characteristics of wings with arbitrary planform, camber, and twist at low to moderate angles of attack. They provide more accurate results than lifting line theory for low aspect ratio wings and complex geometries.}{catAD}
\dictentry{Load Factor}{The ratio of the aerodynamic lift force acting on an aircraft to its weight, expressed in units of gravitational acceleration (g). A load factor of 1 g corresponds to steady, level flight; positive load factors greater than 1 g occur in turns, pull-ups, and gusts, while negative load factors occur in push-overs. Load factor is a fundamental parameter in structural design, flight envelope definition, and maneuver performance assessment.}{catAD}
\dictentry{Longitudinal Stability}{The tendency of an aircraft to resist and recover from disturbances in pitch, such that an increase in angle of attack produces a nose-down restoring moment and a decrease produces a nose-up restoring moment. Longitudinal static stability requires the center of gravity to be ahead of the neutral point, and is provided by the horizontal tail, canard, or flying wing reflex camber. The static margin quantifies the degree of longitudinal stability.}{catAD}
\dictentry{Low Speed Tunnel}{A wind tunnel designed for testing at subsonic speeds well below the compressibility threshold, typically below Mach 0.3, where the flow can be treated as incompressible. Low speed tunnels are the most common type of wind tunnel, used for general aerodynamic research, lift and drag measurements, and flow visualization on a wide range of aircraft, automotive, and civil engineering models. They offer long run times, easy model access, and relatively low operating costs.}{catAD}
\dictentry{Mach Angle}{The angle between the direction of a supersonic flow and a Mach wave (weak disturbance wave) emanating from a point source, given by the relation mu = arcsin(1/M), where M is the local Mach number. The Mach angle defines the boundary of the Mach cone within which the effects of a disturbance in supersonic flow are confined. As the Mach number increases, the Mach angle decreases and the Mach cone becomes more slender.}{catAD}
\dictentry{Mach Cone}{The cone-shaped wavefront that forms around an object moving at supersonic speed, within which all disturbance effects are confined. The half-angle of the Mach cone is the Mach angle, which decreases with increasing Mach number. The region inside the cone is called the zone of action where the flow is influenced by the body, while the region outside is the zone of silence where the flow is unaffected.}{catAD}
\dictentry{Mach Number}{The ratio of the local flow velocity to the local speed of sound in the fluid, named after physicist Ernst Mach. It is the fundamental parameter characterizing compressibility effects in gas dynamics, with M less than 1 indicating subsonic flow, M approximately 1 transonic, 1 to 5 supersonic, and greater than 5 hypersonic. The Mach number determines the character of shock waves, expansion fans, and the overall aerodynamic behavior of a body.}{catAD}
\dictentry{Mach Tuck}{A progressive nose-down pitching moment that develops on an aircraft as it accelerates through the transonic speed regime, caused by the rearward movement of the shock waves and the associated changes in the pressure distribution over the wing and tail surfaces. Mach tuck increases with speed and can overwhelm the trim capability of the elevator if not compensated. It is counteracted by trim tabs, all-moving tails (stabilators), or automatic trim systems.}{catAD}
\dictentry{Mach Wave}{A weak, isentropic disturbance wave that propagates at the Mach angle in supersonic flow, carrying infinitesimally small pressure changes without entropy increase. Unlike shock waves, Mach waves do not alter the flow properties appreciably and represent the limit of influence of a tiny disturbance in supersonic flow. An infinite number of Mach waves form the Mach cone surface that bounds the region of influence of a supersonic body.}{catAD}
\dictentry{Magnetic Suspension}{A wind tunnel model support system that uses electromagnetic forces to levitate and position a model in the test section without any physical struts, arms, or wires, completely eliminating support interference. The model contains permanent magnets or ferromagnetic elements that interact with externally controlled electromagnets to maintain position and orientation. Magnetic suspension enables interference-free measurements but is limited to small models and requires sophisticated control systems.}{catAD}
\dictentry{Maneuverability}{The capacity of an aircraft to change its flight path, speed, and attitude rapidly and precisely in response to pilot commands or automated control inputs. Maneuverability depends on the thrust-to-weight ratio, wing loading, control surface effectiveness, structural load limits, and the aircraft's rotational inertia about each axis. It is a primary performance metric for fighter aircraft and is characterized by parameters such as turn rate, roll rate, and specific excess power.}{catAD}
\dictentry{Maneuvering Load}{The aerodynamic loads imposed on an aircraft structure during pilot-commanded maneuvers such as turns, pull-ups, push-overs, and rolls, expressed in terms of load factor (g). Maneuvering loads are a primary design driver for the wing, fuselage, and control surfaces, and are specified by certification regulations for both limit and ultimate load conditions. The V-n diagram defines the combination of airspeed and load factor that constitutes the flight maneuver envelope.}{catAD}
\dictentry{Manometer}{A device for measuring pressure differences by balancing the weight of a fluid column against the applied pressure difference. In wind tunnel testing, manometers (liquid-column or electronic) are used to measure static pressure at model taps, test section pressure, and tunnel reference pressures. Multi-tube manometers and electronic scanning manometers allow simultaneous measurement of dozens of pressure ports.}{catAD}

\dictentry{Mean Aerodynamic Chord}{The chord length of a hypothetical rectangular wing that would produce the same pitching moment characteristics as the actual wing, calculated as a weighted average of the local chord along the span with the local lift distribution as the weighting function. The MAC is the reference length for moment coefficients, stability derivatives, and center of gravity limits. It is typically shorter than the geometric mean chord for tapered, swept wings.}{catAD}
\dictentry{Mean Line}{The camber line of an airfoil, representing the locus of points equidistant between the upper and lower surfaces, which defines the basic lift-producing curvature of the section. In thin airfoil theory, the airfoil is idealized as its mean line, and the boundary condition of flow tangency is applied on this line rather than on the actual surfaces. The shape of the mean line determines the zero-lift angle and the pitching moment coefficient.}{catAD}
\dictentry{Mesh}{The collection of cells, elements, or nodes that discretize a continuous spatial domain into finite regions for numerical computation in CFD, analogous to a grid but often used to emphasize the connectivity and topology of the discrete elements. Mesh generation is often the most time-consuming step in a CFD workflow, requiring careful attention to cell quality, resolution in regions of high gradients, and smooth transitions between fine and coarse regions. Mesh refinement studies are essential to ensure grid-independent solutions.}{catAD}
\dictentry{Method of Characteristics}{A mathematical technique for solving partial differential equations, particularly hyperbolic equations governing supersonic flow, by reducing them to ordinary differential equations along characteristic lines in the flow field. The method tracks the propagation of disturbances along Mach lines and is used to design supersonic nozzle contours, analyze oblique shock and expansion fan interactions, and compute two-dimensional and axisymmetric supersonic flows. It provides exact solutions for many practical supersonic flow problems.}{catAD}
\dictentry{Moment Coefficient}{A dimensionless parameter that quantifies the aerodynamic moment about a reference point relative to the dynamic pressure, reference area, and reference length (typically mean aerodynamic chord). Cm = M / (q * S * c). The pitching moment coefficient is critical for longitudinal stability and trim analysis, with a negative slope dCm/dCL required for static stability.}{catAD}

\dictentry{Momentum Equation}{Newton's second law applied to fluid motion, stating that the net force acting on a fluid element equals the rate of change of its momentum. In integral form, it relates the sum of surface and body forces to the time rate of change of momentum within a control volume plus the net flux of momentum across its boundaries. The momentum equation is fundamental in aerodynamic force analysis and propulsion system design.}{catAD}
\dictentry{Momentum Thickness}{A boundary layer parameter representing the thickness of a freestream layer that would carry the same momentum deficit as the actual boundary layer, quantifying the loss of momentum due to viscous retardation near the surface. Together with displacement thickness and shape factor, momentum thickness is used to characterize boundary layer state, predict transition, and calculate skin friction drag. It is defined as the integral of the velocity defect weighted by the local velocity across the boundary layer.}{catAD}
\dictentry{Morphing Wing}{A wing structure capable of changing its shape in flight through active mechanisms, smart materials, or compliant structures to optimize aerodynamic performance across multiple flight conditions. Morphing can modify camber, sweep, aspect ratio, span, or surface topology. Potential benefits include improved efficiency, control authority, and mission adaptability, but challenges include weight, complexity, power, and reliability.}{catAD}

\dictentry{Natural Laminar Flow}{A wing and airfoil design approach that shapes the pressure distribution to maintain a favorable (decreasing) pressure gradient over as much of the chord as possible, delaying boundary layer transition from laminar to turbulent without requiring active suction. NLF airfoils feature a long, gradually accelerating flow region on the upper surface followed by a short deceleration region near the trailing edge. The technique can reduce skin friction drag by 20-40 percent compared to turbulent flow designs.}{catAD}
\dictentry{Navier-Stokes Equations}{The fundamental partial differential equations governing viscous fluid motion, expressing the conservation of mass, momentum, and energy for a Newtonian fluid with viscosity and thermal conductivity. They are the most complete description of fluid dynamics and, when solved exactly, predict all phenomena from laminar pipe flow to turbulent aerodynamic flows. Due to their nonlinearity, analytical solutions exist only for very simple cases, and practical solutions require numerical methods.}{catAD}
\dictentry{Neutral Point}{The longitudinal position on an aircraft where the pitching moment coefficient is independent of angle of attack, meaning that a change in angle of attack produces no change in pitching moment. The neutral point is the aft-most position of the center of gravity for which the aircraft remains statically stable, and the distance between the CG and the neutral point defines the static margin. For static stability, the CG must be forward of the neutral point.}{catAD}
\dictentry{Noise Certification}{The regulatory process by which aircraft are certified to meet maximum allowable noise levels established by aviation authorities such as ICAO, FAA, and EASA. Noise certification defines measurement procedures, flight test conditions, and noise limits for takeoff, approach, and sideline.}{catAD}
\dictentry{Noise Reduction}{Techniques and technologies employed to reduce aircraft noise, including engine nacelle treatments, chevron nozzles, acoustic liners, shielding by airframe structures, and operational procedures such as continuous descent approaches.}{catAD}
\dictentry{Normal Force Coefficient}{A dimensionless parameter representing the aerodynamic force component normal to a body's reference plane (typically the chord line for airfoils) relative to dynamic pressure and reference area. Cn = L * cos(alpha) + D * sin(alpha). It is used in missile and rocket aerodynamics where angle of attack varies widely and the force normal to the body axis is the primary control and stability variable.}{catAD}

\dictentry{Normal Shock}{A stationary shock wave in which the flow direction is perpendicular to the shock front, causing an abrupt deceleration from supersonic to subsonic flow with simultaneous increases in pressure, temperature, and $\rho$. Normal shocks occur at the throat of converging-diverging nozzles, at the inlet of supersonic air-breathing engines, and on the surface of blunt bodies at supersonic speeds. The strength of a normal shock increases with the upstream Mach number.}{catAD}
\dictentry{Nose Cap}{The forward-most component of a re-entry vehicle or hypersonic aircraft, exposed to the highest stagnation temperatures and pressures. The nose cap is typically made of reinforced carbon-carbon or an ablative material and must withstand extreme thermal and mechanical loads.}{catAD}

\dictentry{Nusselt Number}{A dimensionless parameter defined as the ratio of convective to conductive heat transfer across a boundary, $Nu = hL / k$. It characterizes the effectiveness of convection relative to pure conduction and is a primary output of convective heat transfer correlations. In aerodynamic heating, Nusselt number distributions determine surface heat flux and thermal protection system requirements.}{catAD}

\dictentry{Oblique Shock}{A shock wave that forms at an angle to the direction of supersonic flow, typically generated by a concave corner, a wedge, or a cone in supersonic freestream. Across an oblique shock, the flow is decelerated and deflected through a finite angle, with increases in pressure, temperature, and entropy that are less severe than those across a normal shock of the same upstream Mach number. Oblique shocks are the primary wave structure in supersonic aerodynamics.}{catAD}
\dictentry{Ogive}{A pointed, smoothly curved shape commonly used for the nose cones of rockets, missiles, and projectiles, defined mathematically as the curve formed by the intersection of a cylinder with a sphere offset from the cylinder axis. Ogive nose shapes provide a favorable balance between low wave drag at supersonic speeds, internal volume for payload, and structural simplicity. Tangent ogives blend smoothly into the cylindrical body, while secant ogives have a distinct junction.}{catAD}
\dictentry{Open Circuit Tunnel}{A wind tunnel configuration in which the test air is drawn from the atmosphere, passed through the test section, and exhausted back to the atmosphere without being recirculated. Open circuit tunnels are simpler and less expensive to construct than closed circuit designs but consume more power since the air is not recycled. They are susceptible to atmospheric disturbances and temperature variations, and produce higher noise levels due to uncontained exhaust.}{catAD}
\dictentry{Optical Strain Measurement}{A non-contact technique that uses cameras, lasers, and digital image correlation (DIC) or moire interferometry to measure surface deformation and strain fields on wind tunnel models or full-scale structures under load. By tracking the displacement of a stochastic speckle pattern or grating applied to the surface, full-field strain maps can be obtained without the mass loading or wiring complications of traditional strain gauges. It is particularly useful for composite structures and complex geometries.}{catAD}
\dictentry{Oscillatory Testing}{Wind tunnel testing in which the model is forced to oscillate in pitch, roll, or yaw at controlled frequencies and amplitudes while the unsteady aerodynamic forces and moments are measured. Oscillatory testing provides the frequency-dependent aerodynamic derivatives (such as pitch damping and roll stiffness) that are essential for predicting flutter boundaries and flight dynamics behavior. The results validate unsteady aerodynamic theories and computational predictions.}{catAD}
\dictentry{Oswald Span Efficiency}{A dimensionless factor (typically 0.7 to 0.9 for practical aircraft) that relates the actual induced drag of a wing to the induced drag of an ideal elliptical wing producing the same total lift and having the same span. Named after W.F. Oswald, it accounts for the non-elliptical lift distribution and parasitic drag interactions of real wings. The efficiency factor is used in preliminary design to estimate induced drag from the lift coefficient, aspect ratio, and wing geometry.}{catAD}
\dictentry{Panel Method}{A numerical technique for solving potential (inviscid, irrotational, incompressible) flow problems by distributing singularities (sources, vortices, or doublets) on panels that approximate the body surface, then solving for the singularity strengths that satisfy the boundary condition of flow tangency. Panel methods are computationally efficient and widely used for preliminary aerodynamic design, providing lift, pressure distribution, and velocity field predictions for complex two-dimensional and three-dimensional body shapes.}{catAD}
\dictentry{Parasite Drag}{The total drag component that is not associated with lift production, consisting of skin friction drag (from viscous shear on wetted surfaces) and form or pressure drag (from the pressure difference between forward and aft faces of the body). Parasite drag is the dominant drag component at high speeds and is minimized through streamlining, smooth surface finishes, and careful attention to component junctions and protuberances. It scales with the square of velocity.}{catAD}
\dictentry{Particle Image Velocimetry}{A whole-field optical flow measurement technique in which the flow is seeded with small tracer particles illuminated by a laser sheet, and two successive images captured by a high-speed camera are cross-correlated in small interrogation windows to determine the local displacement and hence velocity of the particles. PIV provides instantaneous velocity vector maps across a plane or volume, revealing flow structures such as vortices, shear layers, and separation regions. It has become a standard tool in experimental aerodynamics.}{catAD}
\dictentry{Passive Flow Control}{A method of modifying aerodynamic flow fields using fixed geometric features such as vortex generators, surface roughness, ribs, cavities, or shaping that require no energy input, sensors, or active modulation. Passive flow control is simpler, lighter, and more reliable than active control but cannot adapt to changing conditions. Examples include vortex generators to delay separation, serrated trailing edges for noise reduction, and boundary layer trips.}{catAD}

\dictentry{Pathline}{The trajectory or path traced by an individual fluid particle as it moves through a flow field over a period of time, recorded by following that particle from its initial position. Pathlines are relevant for unsteady flows and differ from streamlines (instantaneous velocity direction) and streaklines (locus of particles passing through a fixed point) unless the flow is steady, in which case all three coincide. Understanding pathlines is important for visualizing mixing and pollutant dispersion.}{catAD}
\dictentry{Peak Heating}{The maximum instantaneous heating rate experienced by a spacecraft during atmospheric re-entry, occurring when the combination of velocity and atmospheric $\rho$ produces the highest convective and radiative heat flux. Peak heating determines the required thickness and material of the thermal protection system.}{catAD}
\dictentry{Phase Doppler Anemometry}{An optical measurement technique that simultaneously determines the size and velocity of individual particles or droplets in a flow by analyzing the phase shift between scattered laser signals received at multiple detectors. The phase difference between detectors is proportional to particle diameter, while the Doppler frequency gives velocity. PDA is widely used in spray characterization, combustion research, and multiphase flow studies.}{catAD}
\dictentry{Phugoid}{A long-period, lightly damped longitudinal oscillation of an aircraft in which kinetic energy (airspeed) and potential energy (altitude) are exchanged while the angle of attack remains approximately constant. The phugoid period typically ranges from 30 seconds for light aircraft to several minutes for large transports, and the motion appears as a slow undulating climb and descent. While usually not dangerous, an undamped phugoid can be uncomfortable and requires pilot or autopilot intervention to suppress.}{catAD}
\dictentry{Pitching Moment}{The aerodynamic moment about the lateral (pitch) axis of an aircraft, tending to rotate the nose up or down, and one of the three fundamental moments along with rolling and yawing moments. The pitching moment coefficient and its variation with angle of attack determine the longitudinal static stability and trim characteristics of the aircraft. A nose-up pitching moment is conventionally defined as positive.}{catAD}
\dictentry{Pitot Tube}{A slender, forward-facing probe used to measure the stagnation (total) pressure of a flow by bringing the air to rest isentropically at the probe tip. Combined with a static pressure port, the pitot tube provides the dynamic pressure, from which the flow velocity can be calculated using Bernoulli's equation. Pitot-static probes are fundamental instruments in flight testing, wind tunnel operation, and aircraft airspeed measurement systems.}{catAD}
\dictentry{Prandtl Number}{A dimensionless parameter defined as the ratio of momentum diffusivity (kinematic viscosity) to thermal diffusivity, representing the relative thickness of the velocity boundary layer to the thermal boundary layer. For air at standard conditions, the Prandtl number is approximately 0.71, meaning that momentum and heat diffuse at similar rates. It is a key parameter in convective heat transfer correlations and boundary layer analysis.}{catAD}
\dictentry{Prandtl-Meyer Expansion}{The isentropic expansion process by which a supersonic flow accelerates and turns around a convex corner through a continuous fan of Mach waves, with the flow deflection angle determined by the Prandtl-Meyer function. Across the expansion fan, the Mach number increases while pressure, temperature, and $\rho$ decrease without any entropy change. The process is the exact opposite of an oblique shock and is reversible.}{catAD}
\dictentry{Prandtl-Meyer Function}{A mathematical function of the Mach number that gives the maximum deflection angle a supersonic flow can undergo through a centered expansion fan, denoted as nu(M). It is used to calculate the flow properties after an expansion around a convex corner of known angle. The function increases monotonically with Mach number and is derived from the isentropic flow relations for a perfect gas.}{catAD}
\dictentry{Pressure Coefficient}{A dimensionless number representing the local static pressure on a body surface relative to the freestream static pressure, normalized by the freestream dynamic pressure. A pressure coefficient of zero corresponds to freestream pressure, positive values indicate pressures above freestream (compression), and negative values indicate suction. The distribution of pressure coefficient over a surface is the fundamental output of aerodynamic analysis and determines lift, drag, and moment.}{catAD}
\dictentry{Pressure Drag}{The component of aerodynamic drag arising from the net imbalance of pressure forces acting on the front and rear of a body, caused by flow separation that creates a low-pressure wake behind the body. Also known as form drag or base drag, it depends strongly on the body shape and the degree of flow separation. Streamlining a body to delay separation and reduce the wake size is the primary method for minimizing pressure drag.}{catAD}
\dictentry{Pressure Scanner}{An electronic multi-port pressure measurement device that sequentially samples and digitizes pressure readings from dozens or hundreds of pressure taps connected to a wind tunnel model or test article. Pressure scanners significantly reduce the time and complexity of multi-point pressure surveys compared to individual transducers. They are essential for obtaining detailed surface pressure distributions during wind tunnel testing.}{catAD}
\dictentry{Pressure Sensitive Paint}{A coating applied to wind tunnel models that contains luminescent molecules whose fluorescence intensity or lifetime varies with the local partial pressure of oxygen, and hence with local static pressure. When illuminated with UV light, the paint produces a surface luminescence pattern that can be captured by a camera and converted to a full-field pressure map. PSP provides high spatial resolution pressure data without the need for hundreds of discrete pressure taps.}{catAD}
\dictentry{Pressure Sensor}{A transducer that converts fluid pressure into an electrical signal for measurement, recording, or control. In aerodynamic testing, pressure sensors measure static pressure at model surface taps, total pressure behind shocks, and dynamic pressure in the tunnel flow. Types include piezoresistive, capacitive, and piezoelectric sensors with ranges from sub-mbar to hundreds of bar. Accuracy, frequency response, and temperature stability are critical specifications.}{catAD}

\dictentry{Pressure Transducer}{A pressure sensor that converts pressure into an electrical output signal (typically voltage or current) proportional to the applied pressure. Used interchangeably with pressure sensor in most aerodynamic contexts. High-frequency response transducers are essential for unsteady pressure measurements in flutter, buffet, and aeroacoustic testing.}{catAD}

\dictentry{Propeller}{A rotating airfoil system consisting of two or more blades mounted on a hub that converts engine power into thrust by accelerating a mass of air rearward, generating a pressure increase across the propeller disk. Propellers are highly efficient at low to moderate speeds but lose effectiveness as tip speeds approach Mach 1 due to compressibility effects. Modern propellers use advanced blade shapes, swept tips, and multiple blades to maximize efficiency across the operating range.}{catAD}
\dictentry{Propeller Noise}{The noise generated by rotating propellers due to blade loading, thickness effects, and tip vortices. It consists of tonal components at the blade passage frequency and broadband noise from turbulent boundary layer interactions.}{catAD}
\dictentry{Pulsed Detonation Engine}{An advanced propulsion concept that generates thrust through a series of rapid detonation waves (supersonic combustion waves) rather than the deflagration (subsonic combustion) used in conventional jet engines. Detonation waves achieve a higher thermodynamic efficiency than deflagration by releasing energy at near-constant volume, approaching the ideal detonation cycle. Pulsed detonation engines are under active development for next-generation missiles and launch vehicles.}{catAD}
\dictentry{Radiative Heating}{The component of aerodynamic heating resulting from thermal radiation emitted by the high-temperature shock layer gas, particularly at hypersonic velocities above 10 km/s. Radiative heating becomes dominant at very high re-entry velocities and is a major consideration for Mars return and planetary entry.}{catAD}
\dictentry{Raked Wingtip}{A wingtip with a swept trailing edge that extends the effective wingspan without a vertical fin, reducing induced drag by weakening the tip vortex and distributing the lift more favorably near the tip. Raked wingtips are simpler and lighter than winglets and are used on many modern transport aircraft, including the Boeing 777 and 787. They provide a drag reduction of 3-5 percent compared to a squared-off wingtip.}{catAD}
\dictentry{Ramjet}{An air-breathing jet engine that uses the forward motion of the vehicle to compress incoming air through a diverging inlet, eliminating the need for a mechanical compressor. Ramjets require a minimum flight speed of approximately Mach 2-3 to operate efficiently and cannot produce static thrust, so they require a booster rocket or turbojet for acceleration. They offer high specific impulse at supersonic speeds and are used in some missiles and high-speed target drones.}{catAD}
\dictentry{Range}{The maximum horizontal distance an aircraft can travel on a given fuel or energy supply, achieved by maximizing the distance per unit fuel consumed. Range is optimized by flying at the speed for maximum L/D (propeller) or a specific lift-to-drag ratio (jet). Range depends on initial weight, fuel fraction, specific fuel consumption, and aerodynamic efficiency. Maximum range speed is higher than maximum endurance speed.}{catAD}

\dictentry{Rankine Body}{A three-dimensional, streamlined body shape formed by the superposition of a uniform flow, a source, and a sink in potential flow theory, producing a closed body with a rounded nose and tapered tail. The Rankine body provides an analytical solution for the pressure distribution and drag-free form of a body in ideal flow. It is used as a reference shape in potential flow studies and panel method validation.}{catAD}
\dictentry{RANS}{Reynolds-Averaged Navier-Stokes, a CFD approach in which the instantaneous Navier-Stokes equations are time-averaged and the effects of turbulence on the mean flow are modeled using turbulence closure models such as the k-epsilon or Spalart-Allmaras models. RANS is the most widely used turbulence modeling approach in industrial CFD due to its computational efficiency, though it can be inaccurate for highly separated and unsteady flows. It remains the standard for aircraft aerodynamic design and performance prediction.}{catAD}
\dictentry{Rate of Climb}{The vertical component of aircraft velocity during climbing flight, typically expressed in feet per minute (fpm) or meters per second. Rate of climb is determined by the excess power available above that required for level flight, divided by weight. Maximum rate of climb speed (Vy) is used for optimal altitude gain, while maximum angle of climb speed (Vx) is used for obstacle clearance.}{catAD}

\dictentry{Rate of Turn}{The angular velocity at which an aircraft changes its heading during a turn, typically measured in degrees per second, and determined by the bank angle, true airspeed, and radius of turn. The rate of turn is inversely proportional to airspeed for a given bank angle, so slower aircraft can turn more tightly. Maximum rate of turn is a key performance parameter for fighter aircraft and is limited by structural load factor and aerodynamic capabilities.}{catAD}
\dictentry{Re-entry}{The phase of flight in which a spacecraft or space vehicle returns from orbit or a ballistic trajectory through a planetary atmosphere, experiencing extreme aerodynamic heating, deceleration loads, and communication blackout due to ionized plasma surrounding the vehicle. Successful re-entry requires careful trajectory management and thermal protection systems to dissipate the enormous kinetic energy as heat while keeping the vehicle and crew within survivable limits.}{catAD}
\dictentry{Re-entry Angle}{The angle between the flight path of a re-entering vehicle and the local horizontal at the point of atmospheric entry. A steep re-entry angle increases peak heating rates and deceleration, while a shallow angle extends the heating duration and may cause skip-out.}{catAD}
\dictentry{Re-entry Corridor}{The allowable range of flight path angles within which a re-entering spacecraft must fly to achieve a safe landing at the intended site. Deviations outside the corridor result in excessive heating, structural failure, or landing far from the target.}{catAD}
\dictentry{Re-entry Heating Rate}{The rate at which thermal energy is transferred to a spacecraft surface during atmospheric re-entry, typically expressed in watts per square centimeter (W/cm2). The peak heating rate occurs at the point of maximum deceleration and drives thermal protection system design.}{catAD}
\dictentry{Re-entry Velocity}{The speed at which a spacecraft enters a planetary atmosphere, typically between 7.8 km/s (low Earth orbit) and 11.2 km/s (escape velocity from Earth). The re-entry velocity determines the kinetic energy that must be dissipated as heat during atmospheric passage.}{catAD}
\dictentry{Reattachment}{The process by which a separated boundary layer re-establishes itself as an attached flow downstream of a separation region, typically occurring when the separated shear layer encounters a favorable pressure gradient or when the separated region is bounded. Reattachment creates a secondary boundary layer and is associated with high local heat transfer and pressure recovery. It is observed in shock-boundary layer interactions, behind backward-facing steps, and in cavity flows.}{catAD}
\dictentry{Recirculation}{The reverse flow within a separated boundary layer or wake region where fluid moves upstream relative to the overall flow direction, forming a closed or semi-closed streamline pattern. Recirculation zones occur downstream of separation points and within cavities, and are characterized by low mean velocity, high turbulence, and pressure near the freestream value.}{catAD}

\dictentry{Recirculation Zone}{A region of reversed or circulating flow bounded by the separating and reattaching streamlines, typically found downstream of a bluff body, step, or cavity, or on the leeward side of a high-angle-of-attack wing. The recirculation zone contains vortical structures, has relatively uniform pressure, and strongly influences base drag, heat transfer, and mixing.}{catAD}

\dictentry{Reinforced Carbon Carbon}{A composite material consisting of carbon fibers embedded in a carbon matrix, capable of withstanding temperatures up to 3000F (1650C). Reinforced carbon-carbon is used for the highest-temperature areas of re-entry vehicles, including nose caps and wing leading edges.}{catAD}
\dictentry{Relative Wind}{The direction and velocity of the airflow as experienced by a moving aircraft or airfoil, equal and opposite to the aircraft's velocity through the still air (neglecting wind). The relative wind determines the angle of attack, sideslip angle, and the dynamic pressure acting on the aerodynamic surfaces, and is the fundamental reference for all aerodynamic force and moment calculations. Changes in relative wind direction directly affect lift, drag, and control effectiveness.}{catAD}
\dictentry{Reusable TPS}{A thermal protection system designed to withstand multiple re-entry cycles without replacement, typically using ceramic tiles, fibrous insulation blankets, or metallic panels. Reusable TPS was developed for the Space Shuttle and is essential for affordable reusable launch vehicles.}{catAD}
\dictentry{Reynolds Analogy}{A relationship derived by Osborne Reynolds that connects the skin friction coefficient to the Stanton number (heat transfer coefficient), based on the assumption that the mechanisms of momentum transfer and heat transfer in a turbulent boundary layer are similar. The analogy allows estimation of convective heat transfer rates from friction measurements and vice versa, which is particularly useful when direct heat transfer measurements are difficult. Modified versions of the analogy account for Prandtl number effects.}{catAD}
\dictentry{Reynolds Number}{A dimensionless parameter defined as the ratio of inertial forces to viscous forces in a flow, calculated as the product of velocity, characteristic length, and fluid $\rho$ divided by dynamic viscosity. It determines the flow regime, with low values indicating laminar flow and high values indicating turbulent flow, and governs boundary layer behavior, drag characteristics, and heat transfer. Accurate simulation of Reynolds number is essential for meaningful wind tunnel testing.}{catAD}
\dictentry{Riblet}{Microscopic streamwise grooves or ridges, typically 50-200 micrometers in spacing, applied to aircraft or turbine blade surfaces to reduce turbulent skin friction drag by 5-10 percent. Riblets work by interfering with the cross-stream velocity fluctuations in the near-wall region of the turbulent boundary layer, reducing the momentum transfer to the wall. They have been flight tested on commercial aircraft and are a passive, zero-energy drag reduction technology.}{catAD}
\dictentry{Rogallo Wing}{A flexible, delta-shaped wing design consisting of two conically shaped fabric surfaces held taut by a rigid frame or inflated tubes. Rogallo wings are used in hang gliders, parawings, and some spacecraft recovery systems for their light weight and deployability.}{catAD}
\dictentry{Rolling Moment}{The aerodynamic moment about the longitudinal (roll) axis of an aircraft, tending to rotate one wing up and the other down, and the fundamental mechanism of roll control. Rolling moment is generated primarily by ailerons, elevons, or spoilers, and its magnitude determines the roll rate and responsiveness of the aircraft. A positive rolling moment is conventionally defined as right wing down.}{catAD}
\dictentry{Rotor}{A rotating wing assembly consisting of two or more airfoil-shaped blades mounted on a hub that spins to generate lift and, in some configurations, thrust. In helicopters, the main rotor provides both lift and propulsion, while the tail rotor counteracts torque and provides directional control. Rotor aerodynamics involve complex unsteady phenomena including dynamic stall, blade vortex interaction, and compressibility effects at the advancing blade tips.}{catAD}
\dictentry{Rotor Noise}{The acoustic signature of helicopter main rotors and tail rotors, including blade-vortex interaction noise, high-speed impulsive noise, and broadband turbulence ingestion noise. Rotor noise is a critical consideration for urban air mobility and military operations.}{catAD}
\dictentry{Rudder}{A hinged control surface mounted on the vertical tail of an aircraft, used to control yaw (rotation about the vertical axis). Rudder deflection creates a side force on the tail that generates a yawing moment, and is used for coordinated turns, crosswind correction, and engine-out compensation. At low speeds and high angles of attack, the rudder may be the only effective directional control surface.}{catAD}
\dictentry{Schlieren Photography}{An optical flow visualization technique that makes $\rho$ gradients in transparent media visible by exploiting the refraction of light caused by variations in refractive index associated with $\rho$ changes. Schlieren systems use collimated light and a knife-edge or color filter at the focal point to convert phase shifts into intensity variations, producing detailed images of shock waves, expansion fans, and thermal plumes. It is an indispensable tool in supersonic and hypersonic wind tunnel testing.}{catAD}
\dictentry{Scramjet}{A supersonic combustion ramjet engine in which the airflow remains supersonic throughout the combustion chamber, avoiding the total pressure losses and thermal choking associated with decelerating the flow to subsonic speeds for combustion as in a conventional ramjet. Scramjets can operate at Mach numbers above 5 and are being developed for hypersonic cruise missiles and reusable launch vehicles. The primary challenges include fuel injection, mixing, and stable combustion in the supersonic flow.}{catAD}
\dictentry{Screens}{Fine-mesh grids placed in the settling chamber of a wind tunnel to reduce turbulence intensity and improve flow uniformity. Multiple screens with decreasing mesh size are often used in series to achieve high-quality flow.}{catAD}
\dictentry{Separated Flow}{Flow that has detached from the surface of a body, creating a recirculating wake region between the separated shear layer and the body surface. Separated flow produces large pressure drag, reduces lift, and generates unsteady, turbulent fluctuations that can cause structural fatigue and noise. Understanding and predicting separation is one of the most challenging aspects of aerodynamics and is critical for stall, buffet, and control effectiveness predictions.}{catAD}
\dictentry{Settling Chamber}{The wide-section region of a wind tunnel located between the flow conditioning screens and the contraction cone. It provides a low-velocity environment where turbulence decays and flow uniformity is established before acceleration into the test section.}{catAD}
\dictentry{Shock Diamond}{A repeating pattern of bright diamond-shaped regions visible in the exhaust plume of a supersonic jet engine, formed by the interaction of oblique shock waves and expansion fans that reflect off the jet boundary. The diamonds indicate that the exhaust is either overexpanded or underexpanded relative to the ambient pressure and represent the flow's attempt to equilibrate pressure through wave interactions. The number and intensity of diamonds provide a visual indication of the nozzle operating condition.}{catAD}
\dictentry{Shock Layer}{The region of highly compressed, high-temperature gas between the bow shock wave and the surface of a body traveling at supersonic or hypersonic speed. The shock layer contains the most extreme thermodynamic conditions in the flow field, with temperatures high enough to cause molecular dissociation and ionization at hypersonic speeds. The shock layer properties determine the convective and radiative heating loads on the vehicle surface.}{catAD}
\dictentry{Shock Stall}{A form of aerodynamic stall caused by the interaction between a normal shock wave and the boundary layer on a transonic airfoil, where the shock-induced boundary layer separation causes a sudden loss of lift and increase in drag. Shock stall occurs above the critical Mach number when the shock moves aft with increasing speed and eventually triggers massive separation. It is a key limiting factor in transonic flight and motivates the use of supercritical airfoils.}{catAD}
\dictentry{Shock Tunnel}{A high-enthalpy wind tunnel that uses a strong shock wave propagating through a high-pressure driver section to heat and accelerate the test gas to very high stagnation temperatures and pressures, simulating hypersonic flight conditions. Shock tunnels provide test durations of milliseconds to a few tens of milliseconds, sufficient for force and heat transfer measurements on models. They are used for hypersonic aerodynamics, scramjet testing, and real-gas flow studies.}{catAD}
\dictentry{Shock Wave}{A very thin, nearly discontinuous transition region in a supersonic flow across which pressure, temperature, $\rho$, and entropy increase abruptly while velocity decreases, formed when a supersonic flow encounters a disturbance it cannot propagate ahead of itself. Shock waves are irreversible and dissipative, converting kinetic energy into thermal energy, and are a defining feature of supersonic and hypersonic aerodynamics. Strong shocks produce significant total pressure losses and are the primary source of wave drag.}{catAD}
\dictentry{Shock-Free Airfoil}{An airfoil designed through careful shaping of the pressure distribution so that no shock waves form on the surface at the design Mach number, eliminating wave drag entirely at the design condition. Shock-free airfoils achieve this by ensuring that the local supersonic flow decelerates smoothly through an expansion fan system rather than through a terminating shock. Any deviation from the design condition reintroduces shocks, limiting the practical operating range.}{catAD}
\dictentry{Short Period Mode}{A well-damped, high-frequency longitudinal oscillation of an aircraft in which the angle of attack changes rapidly while the airspeed remains approximately constant, dominated by the pitch stiffness and pitch damping derivatives. The short period mode typically has a period of less than two seconds for most aircraft and is usually well-damped in conventional configurations. Poor short period response can lead to pilot-induced oscillations and is a key consideration in flying qualities assessment.}{catAD}
\dictentry{Sideslip}{The angle between an aircraft's longitudinal axis and the relative wind direction, indicating a lateral component of motion relative to the airflow. Sideslip affects the distribution of aerodynamic forces on the fuselage and vertical tail, and is a critical parameter in lateral-directional stability analysis, crosswind landings, and coordinated turn calculations.}{catAD}
\dictentry{Sink Rate}{The vertical velocity component during descent, typically expressed in feet per minute (fpm). For unpowered aircraft, sink rate at a given airspeed is the inverse of the glide ratio scaled by forward speed. Minimum sink rate speed maximizes time aloft in rising air and is lower than best glide speed. Powered aircraft sink rate is controlled by reducing thrust below the level required for level flight.}{catAD}

\dictentry{Skin Friction}{The tangential shear stress exerted by the flow on a surface due to viscous friction between the slow-moving fluid in the boundary layer and the solid wall. Skin friction is a component of parasite drag and increases with surface roughness, boundary layer turbulence, and Reynolds number. It is the dominant drag component on streamlined bodies at subsonic speeds and is the target of laminar flow control and riblet drag reduction technologies.}{catAD}
\dictentry{Skin Friction Coefficient}{A dimensionless parameter that quantifies the local wall shear stress relative to the freestream dynamic pressure, $C_f = \tau_w / (0.5 \rho V^2)$. It characterizes the viscous drag per unit area and depends on Reynolds number, surface roughness, and boundary layer state (laminar or turbulent). Integration of skin friction coefficient over the wetted area gives total friction drag.}{catAD}

\dictentry{Skin Friction Drag}{The component of aerodynamic drag caused by the tangential viscous shear stresses acting on the wetted surface of a body, proportional to the wetted area and the average skin friction coefficient. It is the dominant drag component on streamlined bodies at subsonic speeds and can account for 40-50 percent of total drag on transport aircraft. Reducing skin friction drag through laminar flow control, surface coatings, or riblets is an active area of research.}{catAD}
\dictentry{Skip Re-entry}{A re-entry trajectory in which the vehicle descends into the upper atmosphere, generates aerodynamic lift to ascend back out of the atmosphere, then re-enters a second time. Skip re-entry extends the range and allows higher entry velocities while managing peak heating.}{catAD}
\dictentry{Slat}{A leading-edge high-lift device that slides or pivots forward and downward from the wing's leading edge to increase the maximum lift coefficient, primarily by creating a slot that energizes the upper-surface boundary layer and delays separation at high angles of attack. Slats extend during takeoff and landing to allow lower approach speeds and shorter field lengths, and are retracted during cruise to maintain a clean leading edge. They work in conjunction with trailing-edge flaps to provide the full high-lift capability.}{catAD}
\dictentry{Slender Body Theory}{An approximate aerodynamic theory for long, thin bodies at small angles of attack that linearizes the governing equations and distributes singularities along the body axis rather than on the surface. Slender body theory provides analytical expressions for the crossflow drag, lift, and moment on slender bodies of revolution, delta wings, and fuselage configurations. It is most accurate for bodies with fineness ratios greater than about 5 and at small angles of attack.}{catAD}
\dictentry{Smoke Visualization}{An experimental flow visualization technique in which smoke or tracer vapor is introduced into the airstream upstream of a model in a wind tunnel to reveal flow patterns, streamlines, separation bubbles, vortex cores, and wake structures. Smoke visualization provides qualitative and sometimes quantitative insight into complex three-dimensional flow phenomena that are difficult to measure with point probes. It is particularly effective for low-speed flows and is used for both research and educational demonstrations.}{catAD}
\dictentry{Solar Powered Flight}{Aircraft propulsion using photovoltaic cells mounted on the wing and fuselage surfaces to convert sunlight into electrical energy for motors and batteries. Solar powered aircraft can achieve sustained flight at high altitudes for extended durations, limited by energy storage and diurnal cycles.}{catAD}
\dictentry{Sonic Boom}{The intense acoustic wave generated by an aircraft flying at supersonic speed, caused by the coalescence of Mach waves into shock waves that propagate to the ground as an N-shaped pressure signature. The sonic boom limits the operational envelope of supersonic aircraft over populated areas and was a major factor in the Concorde's restricted route structure. Research into low-boom designs aims to reshape the shock system to reduce the perceived noise to an acceptable level.}{catAD}
\dictentry{Source Panel Method}{A panel method variant that distributes source singularities on the surface panels of a body to represent the flow displacement effect of the solid boundary, with the source strengths determined by enforcing the flow tangency boundary condition at collocation points. Unlike vortex panel methods, the source panel method does not inherently generate lift and requires an additional Kutta condition or superposition with circulation. It is effective for modeling bodies of revolution and non-lifting components.}{catAD}
\dictentry{Spoiler}{A flat plate or panel hinged at one edge that can be deflected upward from the upper surface of a wing to deliberately disrupt the smooth airflow, destroying lift and increasing drag. Spoilers are used for roll control (differential deflection), speed reduction, and to assist in deceleration after touchdown by dumping lift and transferring weight to the wheels for more effective braking. They are lighter and simpler than conventional ailerons.}{catAD}
\dictentry{Spoileron}{A spoiler that is deflected differentially on each wing to produce a rolling moment by reducing lift on one side while leaving the other wing undisturbed. Unlike ailerons which increase lift on one side and decrease it on the other, spoilerons only destroy lift on the down-going wing, making them less efficient for roll control but simpler mechanically. They are used on some gliders, flying wings, and aircraft with limited aileron authority.}{catAD}
\dictentry{Stabilator}{A one-piece, all-moving horizontal tail surface that pivots about a spanwise axis to provide pitch control, replacing the fixed stabilizer and hinged elevator combination. Stabilators are used on high-speed aircraft because they avoid the control effectiveness loss that conventional elevators experience at transonic speeds due to local shock formation on the fixed portion of the tail. They require stronger actuators since the entire surface must be moved against the aerodynamic hinge moment.}{catAD}
\dictentry{Stagnation Point}{The point on the surface of a body in a flow where the local velocity is zero and the flow has been brought to rest, converting all kinetic energy into pressure and thermal energy. At the stagnation point, the pressure reaches its maximum value (stagnation pressure) and the temperature reaches the stagnation temperature. For blunt bodies in supersonic flow, the stagnation point experiences the highest heating rates and pressure loads.}{catAD}
\dictentry{Stagnation Pressure}{The pressure at a point where the flow has been brought to rest isentropically (without entropy increase), equal to the sum of the static pressure and dynamic pressure in incompressible flow. In compressible flow, the stagnation pressure depends on the Mach number and decreases across shock waves due to entropy generation. It is measured by a pitot tube and is a fundamental parameter in compressible flow calculations and wind tunnel data reduction.}{catAD}
\dictentry{Stall}{The sudden loss of lift and increase in drag that occurs when the angle of attack of a wing exceeds the critical value and the boundary layer separates over a large portion of the upper surface. Stall is caused by the inability of the boundary layer to remain attached against a sufficiently strong adverse pressure gradient, and it results in a dramatic reduction in the lift coefficient. Stall recovery requires reducing the angle of attack below the critical value to reattach the flow.}{catAD}
\dictentry{Starting Vortex}{A clockwise-rotating vortex (when viewed from the right) that is shed from the trailing edge of an airfoil when it first begins to move or when its angle of attack is suddenly changed, as the flow establishes the Kutta condition and develops the bound circulation necessary for lift. The starting vortex is equal and opposite in circulation to the bound vortex that forms on the wing, conserving total circulation in the flow. It eventually dissipates due to viscosity.}{catAD}
\dictentry{Static Margin}{The longitudinal distance between the center of gravity and the neutral point, expressed as a percentage of the mean aerodynamic chord, which quantifies the degree of static longitudinal stability. A positive static margin (CG forward of neutral point) means the aircraft is statically stable, with a larger margin providing stronger restoring moments but also requiring more trim force. Typical values range from 5 to 15 percent MAC for conventional aircraft.}{catAD}
\dictentry{Static Port}{A flush-mounted opening on the aircraft fuselage or probe that measures the ambient static pressure of the undisturbed airflow, one of the two pressure inputs required for the pitot-static airspeed indicator and altimeter. Static ports must be located in areas of undisturbed flow to provide accurate readings, and are often equipped with a static source error correction system to compensate for position errors at different flight conditions.}{catAD}
\dictentry{Static Pressure}{The thermodynamic pressure of a fluid measured by an instrument moving with the flow or by a probe that does not disturb the flow, representing the pressure exerted by the random molecular motion of the gas. Static pressure is distinct from stagnation pressure (which includes dynamic pressure) and is one of the fundamental flow properties used in Bernoulli's equation and compressible flow relations. It varies throughout the flow field in response to velocity changes.}{catAD}
\dictentry{Static Stability}{The initial tendency of an aircraft to return to its original equilibrium position following a small disturbance, assessed by the sign of the restoring moment derivatives. An aircraft with positive static stability in pitch will nose down when disturbed to a higher angle of attack, while one with negative static stability will diverge further. Static stability is a necessary but not sufficient condition for overall dynamic stability and safe flight.}{catAD}
\dictentry{Static Test}{Structural testing in which predetermined loads are applied to an aircraft component or full-scale structure in a fixed configuration, without dynamic or inertial effects, to verify structural strength, stiffness, and load paths. Static tests include ultimate load tests (1.5 times the limit load), fatigue tests, and proof tests, and are required for airworthiness certification. Results are compared with analytical predictions to validate structural design methods.}{catAD}
\dictentry{Static Testing}{Aerodynamic or structural testing conducted under steady-state, time-invariant conditions where the flow properties, model position, and applied loads do not change with time. Static wind tunnel testing provides the baseline lift, drag, and moment characteristics as functions of angle of attack, sideslip, and control surface deflection. The results define the aircraft's trim, stability, and performance envelope and serve as the reference for dynamic testing.}{catAD}
\dictentry{Steady Flow}{A flow in which the velocity, pressure, and all other fluid properties at any fixed point in the flow field remain constant over time, though they may vary from point to point in space. Steady flow is an approximation that holds when the time scales of flow changes are much longer than the observation period, and it greatly simplifies the governing equations by eliminating all time-derivative terms. Most cruise flight aerodynamics can be treated as steady flow.}{catAD}
\dictentry{Sting Mount}{A support strut holding a wind tunnel model from behind, connected through the model's base or a sting cavity. Sting mounts minimize aerodynamic interference compared to strut mounts from above or below, making them preferred for high-speed and aft-section testing. The sting is typically aligned with the model's longitudinal axis and may include internal passages for instrumentation wiring and balance connections.}{catAD}
\dictentry{Strain Gauge}{A sensor that converts mechanical deformation (strain) into a change in electrical resistance, typically using a thin metallic foil pattern bonded to a flexible substrate. Strain gauges are arranged in Wheatstone bridge circuits within force balances to measure aerodynamic loads on wind tunnel models and flight test articles. They provide the high sensitivity and linearity required for accurate multi-component force and moment measurements.}{catAD}
\dictentry{Strake}{A narrow, elongated extension attached to a wing root, fuselage side, or other aerodynamic surface to produce favorable flow interactions. Strakes generate vortices that energize the boundary layer over downstream surfaces, delaying stall and improving high-angle-of-attack handling. They also contribute to directional stability and can reduce interference drag at wing-body junctions.}{catAD}
\dictentry{Strakes}{Sharp-edged longitudinal surfaces designed to generate controlled streamwise vortices that enhance aerodynamic performance at high angles of attack. Strakes are used on fighter aircraft such as the F/A-18 and F-16 to improve lift and stall characteristics, and on missiles to provide additional lift. The vortices they produce remain coherent over a wide range of flight conditions.}{catAD}
\dictentry{Stratospheric Flight}{Flight within the stratosphere (approximately 35,000 to 160,000 feet), characterized by stable atmospheric conditions, low turbulence, and minimal weather. Stratospheric platforms enable long-endurance missions for earth observation, telecommunications, and scientific research.}{catAD}
\dictentry{Streakline}{The locus of all fluid particles that have passed through a fixed point in the flow field at any previous time, visualized by continuously releasing tracer dye or smoke from a stationary source. In steady flow, streaklines coincide with streamlines and pathlines, but in unsteady flow they differ and reveal the history of flow disturbances. Streaklines are commonly used in experimental flow visualization to show vortex formation and wake patterns.}{catAD}
\dictentry{Streamline}{A line in the flow field that is everywhere tangent to the local velocity vector at a given instant, providing an instantaneous snapshot of the flow direction. Streamlines cannot cross one another (except at stagnation points) and their spacing indicates relative flow speed, with closer spacing indicating higher velocity. The collection of streamlines defines the streamline pattern that visualizes the entire flow field at a moment in time.}{catAD}
\dictentry{Strouhal Number}{A dimensionless parameter defined as the ratio of oscillatory frequency times characteristic length to flow velocity, $St = fL / V$. It characterizes the dominant frequency of vortex shedding from bluff bodies, the unsteadiness of separated flows, and the scaling of oscillatory aerodynamic phenomena. Strouhal numbers for circular cylinders are approximately 0.2 over a wide Reynolds number range.}{catAD}

\dictentry{Strut Mount}{A support structure that holds a wind tunnel model from below through the floor or from above through the ceiling of the test section. Strut mounts provide stable, rigid support and allow easy model access and configuration changes between runs. They introduce more flow interference than sting mounts and require careful aerodynamic fairings to minimize their effect on measurements.}{catAD}
\dictentry{Suborbital Flight}{A trajectory that reaches space (above the Karman line at 100 km altitude) but lacks the velocity to achieve orbit, resulting in a ballistic return to Earth. Suborbital flight is used for sounding rockets, crewed space tourism, and hypersonic research.}{catAD}
\dictentry{Subsonic Flow}{Flow in which the local Mach number is less than 1 everywhere, meaning the flow velocity is below the speed of sound and pressure disturbances can propagate both upstream and downstream. In subsonic flow, compressibility effects are negligible below Mach 0.3, and the flow can be treated as incompressible. Most general aviation, commercial transport cruise, and UAV flight operations are conducted in the subsonic regime.}{catAD}
\dictentry{Supercritical Airfoil}{An airfoil designed with a relatively flat upper surface and a highly cambered aft section to delay the onset of drag rise at transonic speeds. By flattening the upper surface, the supercritical airfoil reduces the local flow acceleration and weakens the terminating shock wave, extending the drag rise Mach number. It allows thicker wings with more internal volume while maintaining efficient transonic cruise performance.}{catAD}
\dictentry{Supersonic Flow}{Flow in which the local Mach number exceeds 1, characterized by the formation of shock waves and expansion fans that fundamentally change the aerodynamic behavior. In supersonic flow, pressure disturbances cannot propagate upstream and are confined within Mach cones, leading to wave drag as a new dominant force component. The governing equations change from elliptic to hyperbolic type, requiring specialized analysis methods.}{catAD}
\dictentry{Supersonic Wind Tunnel}{A wind tunnel designed to operate at Mach numbers above 1, using a converging-diverging nozzle to accelerate the test section flow to supersonic speeds. The test section typically has a fixed or adjustable geometry to achieve a range of Mach numbers, and the tunnel must handle the high power requirements and pressure ratios needed to maintain supersonic flow. Schlieren or shadowgraph optics are standard for visualizing the shock waves that form around models.}{catAD}

\dictentry{Sweep Angle}{The angle measured between the wing quarter-chord line and a line perpendicular to the aircraft's longitudinal axis, indicating how far the wing is swept back. Sweep angle is a primary parameter influencing transonic and supersonic aerodynamic performance, as it reduces the effective Mach number normal to the leading edge. Higher sweep angles delay drag rise but increase structural weight and reduce low-speed lift capability.}{catAD}
\dictentry{Sweep Theory}{A simplified method for estimating the aerodynamic characteristics of swept wings by applying two-dimensional airfoil data resolved along the direction normal to the leading edge. By decomposing the freestream velocity into components normal and parallel to the sweep line, the theory accounts for the reduced effective Mach number seen by each wing section. Sweep theory provides rapid preliminary estimates of lift, drag, and pitching moment for swept wing configurations.}{catAD}
\dictentry{Swept Wing}{A wing planform in which the wing is angled rearward from root to tip, with the leading edge and trailing edge both swept back relative to the lateral axis. Sweeping the wing increases the critical Mach number by reducing the effective Mach component normal to the leading edge, delaying the onset of compressibility drag rise. Most modern high-subsonic and supersonic aircraft use swept wings to achieve efficient cruise performance.}{catAD}
\dictentry{T-Tail}{A tail empennage configuration in which the horizontal stabilizer is mounted at the top of the vertical fin, forming a T shape when viewed from the front. The T-tail places the horizontal surface outside the wing wake and jet exhaust, improving pitch control effectiveness at low speeds. It also provides beneficial endplate effects that can reduce the required vertical tail size, though it introduces structural weight and flutter concerns.}{catAD}
\dictentry{Taileron}{A combined stabilator and aileron control surface used on fighter aircraft such as the F-16 and MiG-29, where the entire left and right horizontal tails deflect symmetrically for pitch control and differentially for roll control. Tailerons eliminate the need for separate ailerons by using the elevons on each side of the tail independently. They provide powerful pitch and roll authority at all flight speeds.}{catAD}
\dictentry{Takeoff Distance}{The horizontal distance required for an aircraft to accelerate from rest to a defined screen height (typically 35 or 50 feet) above the runway, including the ground roll and initial climb segment. Takeoff distance depends on weight, $\rho$ altitude, wind, runway gradient, thrust, and pilot technique. Certification regulations require balanced field length for multi-engine aircraft, where the distance to accelerate to V1 and stop equals the distance to continue takeoff and clear the screen height after engine failure at V1.}{catAD}

\dictentry{Taper Ratio}{The ratio of the tip chord to the root chord of a wing, expressed as a dimensionless number between 0 and 1. Taper ratio influences the spanwise lift distribution, structural weight, and manufacturing complexity, with moderate taper (0.3-0.5) providing a good compromise between aerodynamic efficiency and structural practicality. An elliptically tapered wing approximates the ideal elliptical lift distribution for minimum induced drag.}{catAD}
\dictentry{Temperature Sensitive Paint}{A coating applied to wind tunnel models that changes color or luminescence in response to surface temperature variations, enabling full-field temperature mapping. TSP allows researchers to identify boundary layer transition by detecting the temperature rise associated with turbulent flow, locate shock impingement points, and measure heat transfer distributions. It provides high spatial resolution without requiring discrete sensor installations.}{catAD}
\dictentry{Test Section}{The portion of a wind tunnel where the model is placed and aerodynamic measurements are conducted, designed to provide a uniform, controlled flow of known velocity and direction. Test sections may be open or closed, with solid, slotted, or perforated walls depending on the speed range and testing requirements. The cross-sectional shape and size determine the maximum model dimensions and blockage ratio.}{catAD}
\dictentry{Thermal Barrier Coating}{A thin layer of ceramic material applied to metallic components to reduce heat transfer and protect against oxidation and thermal fatigue. Thermal barrier coatings are used on turbine blades, combustion chambers, and exhaust nozzles in gas turbine engines.}{catAD}
\dictentry{Thermal Protection Material}{Materials specifically engineered to withstand the extreme heat loads of atmospheric re-entry or hypersonic flight. Thermal protection materials include ablative composites (carbon-phenolic), reusable ceramics (tiles), and advanced carbon-carbon composites, each selected for specific temperature ranges and mission profiles.}{catAD}
\dictentry{Thermal Protection System}{The engineered combination of materials and structures that protects a spacecraft from the extreme aerodynamic heating encountered during atmospheric re-entry or hypersonic flight. Thermal protection systems must withstand temperatures ranging from hundreds to thousands of degrees while maintaining structural integrity and shape. They can be ablative, consuming material to dissipate heat, or reusable, using ceramic tiles or metallic panels that radiate heat away.}{catAD}
\dictentry{Thermocouple}{A temperature sensor consisting of two dissimilar metal wires joined at one end, producing a voltage proportional to the temperature difference between the measuring junction and a reference junction. In aerodynamic testing, thermocouples are embedded in wind tunnel models and thermal protection materials to measure surface temperatures and heat fluxes. Their small size, ruggedness, and wide temperature range make them ideal for harsh aerodynamic environments.}{catAD}
\dictentry{Thickness Ratio}{The ratio of the maximum thickness of an airfoil section to its chord length, expressed as a percentage, defining the relative slenderness of the profile. A higher thickness ratio provides more internal volume for fuel and structure but increases drag, especially at transonic and supersonic speeds. Typical thickness ratios range from 12-15 percent for subsonic transports to 3-6 percent for supersonic fighters.}{catAD}
\dictentry{Thin Airfoil Theory}{A linearized potential flow theory developed by Munk and Glauert that predicts the lift and pitching moment characteristics of a thin, slightly cambered airfoil at small angles of attack. The theory replaces the airfoil with its camber line, distributes vorticity along the chord, and applies the flow tangency condition to solve for the circulation distribution. It provides analytical expressions for lift slope, zero-lift angle, and moment coefficient.}{catAD}

\dictentry{Thrust Vectoring}{The ability to redirect the exhaust flow of a jet engine to produce forces and moments beyond those generated by conventional aerodynamic control surfaces. By deflecting the exhaust nozzle in pitch or yaw, thrust vectoring provides enhanced maneuverability, post-stall agility, and shorter takeoff distances. It is used on advanced fighters like the F-22 Raptor and Su-35 to achieve supermaneuverability.}{catAD}
\dictentry{Tiles}{Rigid ceramic insulating blocks used as reusable thermal protection on the Space Shuttle and other re-entry vehicles. Tiles are made of high-purity silica fibers and are coated with a borosilicate glass coating for waterproofing and emissivity control.}{catAD}
\dictentry{Tip Fence}{A small vertical or angled surface mounted at the wingtip, functioning as a simplified endplate to reduce induced drag by partially blocking the crossflow from the lower to the upper surface. Tip fences are lighter and simpler than winglets while providing a modest improvement in effective aspect ratio. They are found on some business jets, sailplanes, and regional aircraft as a cost-effective drag reduction measure.}{catAD}
\dictentry{Total Pressure}{The sum of static pressure and dynamic pressure in a flow, representing the pressure that would be measured if the flow were brought to rest isentropically at a stagnation point. Total pressure remains constant in isentropic flow but decreases across shock waves due to entropy generation, making it a key indicator of losses. It is measured by a pitot tube and used in compressible flow analysis and engine performance assessment.}{catAD}
\dictentry{Trailing Edge}{The aft-most edge of an airfoil, wing, or control surface where the upper and lower surface flows rejoin and the boundary layer leaves the surface. The trailing edge geometry affects the Kutta condition, base drag, and control surface hinge moments. Sharp trailing edges are preferred for aerodynamic efficiency, while blunt trailing edges may be used for structural or manufacturing reasons.}{catAD}
\dictentry{Trailing Vortex}{The powerful, counter-rotating spiral vortex that rolls up from the wingtip region behind an aircraft as a result of the pressure difference between the upper and lower surfaces. Trailing vortices constitute the wake vorticity that is responsible for induced drag and pose a hazard to following aircraft due to the strong rotational velocities they contain. They persist for several minutes behind large transport aircraft and govern aircraft separation standards.}{catAD}
\dictentry{Transitional Flow}{The intermediate flow regime in which the boundary layer undergoes transition from a stable laminar state to a fully turbulent state, characterized by the growth and breakdown of Tollmien-Schlichting waves. Transition typically occurs over a finite streamwise distance and can be triggered by surface roughness, freestream turbulence, pressure gradients, or acoustic disturbances. The transition location profoundly affects skin friction drag and heat transfer.}{catAD}
\dictentry{Transonic Flow}{The flow regime between approximately Mach 0.8 and 1.2 where both subsonic and supersonic regions coexist in the flow field, typically with supersonic pockets on the upper surface of wings terminated by shock waves. Transonic flow is characterized by the onset of wave drag, shock-induced separation, and buffet, making it a critical design regime for commercial transport aircraft. The mixed nature of the flow challenges both computational and experimental methods.}{catAD}
\dictentry{Trim Drag}{The additional drag generated when control surfaces, typically elevators or stabilators, are deflected from their neutral positions to maintain trimmed flight at a given airspeed and center of gravity location. Trim drag results from the change in the aircraft's effective camber and the induced drag from the tail lift required to balance the pitching moment. It is minimized by careful CG management and aerodynamic tailoring of the wing-tail interaction.}{catAD}
\dictentry{Turbine}{A rotating turbomachinery component that extracts energy from a high-temperature, high-pressure gas stream by expanding it through alternating rows of stationary nozzle vanes and rotating rotor blades. In a gas turbine engine, the turbine drives the compressor and fan, converting the thermal energy of combustion into shaft power. Turbine performance is characterized by expansion ratio, efficiency, and the ability to withstand extreme temperatures.}{catAD}
\dictentry{Turbofan}{A ducted jet engine configuration in which a large front fan driven by the turbine accelerates a large mass of air through the bypass duct while the core engine generates a separate exhaust jet. Most of the thrust in a high-bypass-ratio turbofan comes from the fan, providing superior fuel efficiency and lower noise compared to turbojets. Turbofans are the standard propulsion system for virtually all modern commercial transport aircraft.}{catAD}
\dictentry{Turbojet}{A gas turbine engine that draws in air, compresses it in an axial or centrifugal compressor, mixes it with fuel and combusts it, then expands the hot gas through a turbine and exhaust nozzle to generate thrust. Turbojets have high specific thrust and are efficient at high speeds but consume more fuel than turbofans at subsonic cruise. They are used primarily on supersonic aircraft, missiles, and early-generation jet fighters.}{catAD}
\dictentry{Turbulence Intensity}{A dimensionless measure of the fluctuation level in a wind tunnel flow, defined as the root-mean-square of the velocity fluctuations divided by the mean flow velocity. Low turbulence intensity, typically below 0.1 percent, is required for accurate boundary layer transition studies and aerodynamics research. Screens, honeycombs, and high contraction ratios are used to reduce turbulence intensity in wind tunnels.}{catAD}
\dictentry{Turbulence Model}{A mathematical closure model that approximates the Reynolds stresses and turbulent heat fluxes in the Reynolds-Averaged Navier-Stokes (RANS) equations, providing a practical way to compute turbulent flows without resolving all scales. Common turbulence models include the Spalart-Allmaras one-equation model, the k-epsilon and k-omega two-equation models, and the SST blending model. The choice of model depends on the flow type, available computational resources, and accuracy requirements.}{catAD}
\dictentry{Turbulent Flow}{Chaotic, irregular fluid motion characterized by random fluctuations in velocity and pressure at multiple scales, with efficient mixing, high momentum transfer, and increased dissipation compared to laminar flow. Turbulent boundary layers are thicker, more resistant to separation, and produce significantly higher skin friction than laminar ones. Most practical aerodynamic flows at flight Reynolds numbers are turbulent.}{catAD}
\dictentry{Turbulent Wake}{The region of highly turbulent, velocity-deficit flow behind a body where the boundary layers and separated shear layers mix and evolve downstream. The turbulent wake contains the vorticity shed from the body, has lower momentum than the freestream, and expands due to entrainment and turbulent diffusion. Wake properties determine base drag and influence downstream surfaces such as tail surfaces and following aircraft.}{catAD}

\dictentry{Turning Radius}{The radius of the circular flight path traced by an aircraft during a coordinated turn at a constant airspeed and bank angle. The turning radius decreases with increasing bank angle and decreasing airspeed, and is also affected by the load factor and wing loading. Minimum turning radius is a key performance measure for fighter aircraft.}{catAD}
\dictentry{Turning Vanes}{Aerodynamic surfaces placed in duct bends or wind tunnel corners to redirect airflow efficiently. They prevent flow separation and reduce total pressure losses by guiding the flow around turns.}{catAD}
\dictentry{Twin Tail}{An empennage configuration with two separate vertical fins mounted on the aft fuselage or at the ends of the horizontal stabilizer, providing directional stability and control redundancy. Twin tails are common on multi-engine aircraft, carrier-based aircraft, and some fighter designs where single vertical tail height is limited by hangar clearance. They also place the vertical surfaces in undisturbed flow outside the wing wake.}{catAD}
\dictentry{Unsteady Flow}{A flow condition in which the velocity, pressure, $\rho$, and other fluid properties at a given point vary with time, requiring the inclusion of time-derivative terms in the governing equations. Unsteady aerodynamic phenomena include flutter, dynamic stall, gust response, vortex shedding, and maneuvering flight. Analysis of unsteady flow is essential for predicting aeroelastic stability, aircraft response to atmospheric disturbances, and rotorcraft aerodynamics.}{catAD}
\dictentry{Upwash}{The upward component of flow velocity induced ahead of a lifting wing or airfoil, caused by the circulation around the wing that deflects the incoming flow upward before it reaches the leading edge. Upwash affects the effective angle of attack of surfaces mounted ahead of the wing, such as canards and the tail of a T-tail configuration. It also influences the flow over wings in biplane or close-coupled canard arrangements.}{catAD}
\dictentry{V-n Diagram}{A graphical representation of the structural flight envelope plotting airspeed on the horizontal axis against load factor (g) on the vertical axis, defining the permissible combinations for safe flight. The diagram shows maneuver load limits, gust load lines, stall boundaries, and dive speeds as specified by certification regulations. All points within the enclosed area are structurally safe for the aircraft.}{catAD}
\dictentry{V-Tail}{A tail empennage configuration consisting of two angled surfaces meeting in a V shape, each combining the functions of a vertical fin and a horizontal stabilizer. The V-tail reduces the number of surfaces and intersections, lowering drag and radar cross-section, but requires complex mixing of control inputs to achieve independent yaw and pitch control. It is used on aircraft such as the Beechcraft Bonanza and Rutan VariEze.}{catAD}
\dictentry{Vapor Cone}{A visible cone of condensed water vapor that forms around an aircraft flying at transonic speeds in humid conditions, caused by the rapid expansion and cooling of air in regions of local supersonic flow. The vapor cone appears near the wing upper surface, fuselage, or tail where the pressure drops below the saturation point. It indicates the presence of supersonic pockets and shock waves even though the aircraft may not be fully supersonic.}{catAD}
\dictentry{Viscosity}{A measure of a fluid's resistance to shear deformation, arising from intermolecular forces and momentum transfer between fluid layers. In aerodynamics, viscosity governs the behavior of the boundary layer, determines skin friction drag, and influences flow separation, transition to turbulence, and heat transfer.}{catAD}
\dictentry{Vortex}{A region in a fluid flow where the flow rotates about a common centerline, characterized by concentrated vorticity and a distinct swirling motion. Vortices form at wingtips (trailing vortices), at leading edges of delta wings (leading-edge vortices), and behind bluff bodies (von Karman vortex street). They play a central role in lift generation, induced drag, and flow control.}{catAD}
\dictentry{Vortex Generator}{A small vane, tab, or fin mounted on the surface of an airfoil, wing, or control surface that produces streamwise vortices to energize the boundary layer and delay flow separation. By mixing high-momentum free-stream air into the low-momentum near-wall flow, vortex generators increase the boundary layer's resistance to adverse pressure gradients. They are used on wings, engine nacelles, and tail surfaces to improve stall characteristics and control effectiveness.}{catAD}
\dictentry{Vortex Lattice Method}{A numerical aerodynamic panel method that models lifting surfaces as a lattice of vortex ring or horseshoe vortex elements distributed on the mean camber surface. VLM solves for the circulation strengths that satisfy the flow tangency condition at control points, providing predictions of lift distribution, induced drag, and pitching moment. It is widely used in preliminary aircraft design for its computational efficiency and reasonable accuracy.}{catAD}
\dictentry{Vortex Shedding}{The periodic, alternating detachment of vortices from the sides of a bluff body placed in a flow, forming the von Karman vortex street in the wake. Vortex shedding produces oscillatory lift and drag forces on the body at a frequency related to the Strouhal number, which can excite structural vibrations if it matches a natural frequency. It is an important consideration in the design of structures exposed to wind.}{catAD}
\dictentry{Wake}{The region of retarded, turbulent, and often separated flow that forms behind a body moving through a fluid, consisting of the momentum deficit region downstream of the boundary layer separation or trailing edge. The wake contains the vorticity shed from the body and is characterized by velocity deficits, turbulent fluctuations, and pressure lower than the freestream. Wake properties determine the base drag and influence downstream surfaces.}{catAD}
\dictentry{Wake Turbulence}{The turbulent velocity field left behind an aircraft by the trailing vortex system and the momentum deficit wake, which poses a hazard to following aircraft due to strong rotational velocities and sudden gusts. Wake turbulence intensity scales with aircraft weight and inversely with wingspan and speed. Separation standards for takeoff and landing are based on wake turbulence categories (heavy, large, small) to ensure following aircraft can safely traverse the wake.}{catAD}

\dictentry{Wall Interference}{The distortion of the flow field around a wind tunnel model caused by the presence of the test section walls, which artificially constrain the streamlines compared to free-air conditions. Wall interference leads to errors in measured forces, moments, and pressures, and must be corrected using theoretical or empirical methods. Perforated or slotted walls reduce interference at transonic speeds by allowing flow through the boundaries.}{catAD}
\dictentry{Wall Shear Stress}{The tangential frictional stress exerted by the fluid on a solid surface at the wall, arising from the velocity gradient at the wall in the boundary layer. Wall shear stress is directly related to skin friction drag and is a key parameter in boundary layer analysis and turbulence modeling. It can be measured directly with floating-element balances or inferred from velocity profile measurements.}{catAD}
\dictentry{Wash-In}{A geometric twist in a wing that increases the local angle of incidence from root to tip, causing the tip sections to operate at a higher angle of attack than the root. Wash-in can be used to tailor the spanwise lift distribution but is generally avoided because it promotes tip stall, which can lead to aileron effectiveness loss and undesirable stall characteristics at the critical moment of landing.}{catAD}
\dictentry{Washout}{A geometric twist in a wing that progressively decreases the local angle of incidence from root to tip, causing the root sections to stall before the tips. Washout improves lateral control effectiveness at the stall and provides more predictable stall behavior by ensuring that the wingtips remain unstalled when the root begins to separate. It is standard design practice on most aircraft to ensure benign stall characteristics.}{catAD}
\dictentry{Wave Drag}{The component of aerodynamic drag caused by the formation of shock waves on a body in transonic or supersonic flow, representing the irreversible conversion of kinetic energy into thermal energy across the shock. Wave drag increases rapidly with Mach number above the drag rise Mach number and is minimized through area ruling, slender body shaping, and swept or supersonic wing designs. It sets the minimum drag achievable at supersonic speeds.}{catAD}
\dictentry{Weber Number}{A dimensionless parameter defined as the ratio of inertial forces to surface tension forces, We = $\rho V^2$ * L / sigma. It characterizes the importance of surface tension effects on free-surface flows, droplet formation, spray atomization, and cavitation inception. In aerodynamics, it is relevant for liquid fuel injection, rain ingestion, and icing phenomena.}{catAD}

\dictentry{Wind Tunnel}{An experimental testing facility in which a controlled, uniform airflow is passed over a scale model or full-scale test article to measure aerodynamic forces, moments, and flow characteristics. Wind tunnels range in size from small educational units to massive facilities capable of testing full-size aircraft, and operate at speeds from subsonic to hypersonic. They remain an essential complement to computational fluid dynamics in aircraft and vehicle development.}{catAD}
\dictentry{Wind Tunnel Fan}{The fan or drive system that generates the airflow in a wind tunnel, powered by electric motors or gas turbines depending on the tunnel size and speed requirements. The fan must overcome the total pressure losses through the tunnel circuit, including those from the test section, diffuser, corners, and flow conditioning elements. Variable-pitch fan blades allow precise control of test section velocity.}{catAD}
\dictentry{Wing}{The primary lifting surface of an aircraft, designed to generate the aerodynamic force of lift that opposes the aircraft's weight and enables flight. Wings are typically airfoil-shaped with cambered upper and lower surfaces, and their geometry includes planform shape, span, chord, sweep, and thickness distribution. The wing absorbs flight loads, houses fuel in many designs, and provides attachment points for engines and landing gear.}{catAD}
\dictentry{Wing Area}{The planform surface area of the wing, defined as the projected area when viewed from above, used as the reference area for nondimensional lift, drag, and moment coefficients. Wing area directly affects the wing loading and thereby the stall speed, maneuverability, and cruise efficiency. It is a fundamental parameter in aircraft sizing and performance calculations.}{catAD}
\dictentry{Wing Fence}{A thin vertical surface mounted chordwise on the upper surface of a swept wing to inhibit spanwise boundary layer flow toward the wingtip. By blocking the crossflow component, wing fences improve aileron effectiveness at high angles of attack, delay tip stall, and provide more predictable stall characteristics. They were a common feature on early swept-wing jet aircraft such as the MiG-15 and Sabre.}{catAD}
\dictentry{Wing Loading}{The ratio of the aircraft's gross weight to its wing area, expressed in pounds per square foot or newtons per square meter, reflecting the average pressure the wing must support in level flight. Low wing loading produces lower stall speeds and better maneuverability, while high wing loading improves cruise efficiency and gust penetration. Wing loading is a primary design parameter that influences the aircraft's performance envelope.}{catAD}
\dictentry{Wing Root}{The structural and geometric junction where the wing attaches to the fuselage, transferring the bending moments, shear forces, and torsional loads from the wing into the airframe. The wing root experiences the highest stresses of any wing station and is typically reinforced with thicker skins, heavier spars, and complex fittings. Root fairings are often added to reduce interference drag at the wing-body junction.}{catAD}
\dictentry{Wing Span}{The total distance measured from one wingtip to the other, perpendicular to the aircraft centerline, determining the aspect ratio and aerodynamic efficiency of the wing. Longer spans reduce induced drag for a given wing area but increase structural weight and bending moments. Wing span is a primary constraint in aircraft design, limited by hangar size, airport gate compatibility, and structural dynamics.}{catAD}
\dictentry{Wing Tip}{The extreme outer end of the wing, where the upper and lower surface flows meet and the trailing vortex is generated. Wingtip geometry significantly influences induced drag and overall wing efficiency, with devices like winglets, tip fences, and raked tips reducing the strength of the tip vortex. The tip also often houses navigation lights, static discharge wicks, and antennae.}{catAD}
\dictentry{Winglet}{A small vertical or angled extension mounted at the wingtip that reduces induced drag by diffusing the tip vortex and recovering some of the energy otherwise lost in the wingtip swirl. Winglets effectively increase the aspect ratio without extending the physical wingspan, improving cruise fuel efficiency by 3-5 percent. They are ubiquitous on modern transport aircraft, business jets, and sailplanes.}{catAD}
\dictentry{Wingtip Vortex}{The powerful, coherent spiral vortex that forms at each wingtip due to the pressure-driven flow of air from the high-pressure lower surface around the tip to the low-pressure upper surface. Wingtip vortices contain much of the kinetic energy of the trailing vortex system and are the primary cause of induced drag on a finite wing. They can persist for miles behind large aircraft and create turbulence hazards for following traffic.}{catAD}
\dictentry{Yawing Moment}{The aerodynamic moment about the vertical (yaw) axis of an aircraft, tending to rotate the nose to the left or right. Yawing moment is generated by the vertical tail, rudder deflection, differential thrust, and asymmetric drag. It is the fundamental moment in directional stability and control analysis, determining the aircraft's ability to maintain coordinated flight and compensate for engine failure.}{catAD}

\dictentry{Zero-Lift Angle of Attack}{The angle of attack at which an airfoil or wing generates zero lift, a fundamental parameter in thin airfoil theory and wing design. For symmetric airfoils the zero-lift angle is zero, while for cambered airfoils it is negative, meaning the airfoil produces lift even at zero geometric angle of attack. The zero-lift angle depends only on the camber line shape and is independent of thickness distribution and Reynolds number in inviscid theory.}{catAD}

\dictentry{Zero-Lift Drag Coefficient}{The drag coefficient of an airfoil, wing, or aircraft at the zero-lift angle of attack, representing the baseline parasitic drag excluding induced drag. It is a fundamental performance parameter that includes skin friction drag, pressure drag from thickness, and interference drag. Minimizing zero-lift drag is essential for achieving high cruise efficiency and maximum range.}{catAD}

\dictentry{Zonal Method}{A computational fluid dynamics approach that partitions the flow domain into distinct zones, each solved with the most appropriate method or level of approximation. Common zonal strategies couple a Navier-Stokes solver near the body with an Euler or potential flow solver in the far field, or use different turbulence models in different regions. Zonal methods optimize computational cost by applying high-fidelity methods only where they are needed for accuracy.}{catAD}

